\documentclass[10pt]{article}
\usepackage[a5paper,margin=11mm]{geometry}
% Общая преамбула для решений и шпаргалки
\usepackage[T2A]{fontenc}
\usepackage[utf8]{inputenc}
\usepackage[russian]{babel}
\usepackage{amsmath,amssymb}
\usepackage{xcolor}
\usepackage[most]{tcolorbox}
\usepackage{enumitem}
\usepackage{needspace}

\setlist{nosep,leftmargin=*}
\setlength{\parindent}{0pt}
\setlength{\parskip}{3pt}

% --- обозначения ---
\newcommand{\dd}{\mathrm{d}}
\newcommand{\avg}[1]{\overline{#1}}
\newcommand{\vp}{v_{\text{в}}}            % наивероятнейшая скорость
\newcommand{\vav}{\bar v}                  % средняя скорость
\newcommand{\erf}{\operatorname{erf}}

% --- задача ---
% \begin{problem}{1.3}{dz}{условие} ... \end{problem}
% второй аргумент: dz — задача из ДЗ, ex — разобранный пример книги,
% exdz — пример, который одновременно входит в ДЗ
\definecolor{dzcol}{RGB}{180,40,40}
\definecolor{excol}{RGB}{40,90,170}
\makeatletter
\newcommand{\tag@dz}{\textcolor{dzcol}{ДЗ}}
\newcommand{\tag@ex}{\textcolor{excol}{пример}}
\newcommand{\tag@exdz}{\textcolor{excol}{пример}\,+\,\textcolor{dzcol}{ДЗ}}
\newcommand{\probtag}[1]{\csname tag@#1\endcsname}
\makeatother
\newtcolorbox{problem}[3]{enhanced,breakable,
  colback=white,colframe=black!55,boxrule=0.5pt,arc=1pt,
  left=4pt,right=4pt,top=3pt,bottom=3pt,before skip=6pt,after skip=6pt,
  fonttitle=\bfseries,coltitle=black,colbacktitle=black!8,
  title={Задача #1\hfill{\normalfont\small\bfseries\probtag{#2}}},
  before upper={\textit{#3}\par\smallskip}}

% блок опорных формул в начале параграфа
\newtcolorbox{basics}{enhanced,breakable,colback=blue!3,colframe=excol!70,
  boxrule=0.5pt,arc=1pt,left=4pt,right=4pt,top=3pt,bottom=3pt,
  fonttitle=\bfseries,title={Опорные формулы}}
\newcommand{\sub}[1]{\par\smallskip\textbf{#1}\par}

\newcommand{\idea}[1]{\textbf{Идея:} #1\par}
\newtcolorbox{answerbox}{colback=yellow!12,colframe=black!40,boxrule=0.4pt,arc=1pt,
  left=3pt,right=3pt,top=1pt,bottom=1pt,before skip=4pt,after skip=2pt}
\newcommand{\answer}[1]{\begin{answerbox}\textbf{Ответ:} #1\end{answerbox}}
\newcommand{\note}[1]{\par{\small\textcolor{black!60}{\textit{Замечание:} #1}}\par}


\title{Решения задач ДЗ и разобранных примеров\\[2pt]
\large§1, §2, §3, §8, §9 — Вишнякова, Макаров и др.\ «Основы термодинамики и статистической физики»}
\date{}

\begin{document}
\maketitle
\vspace{-8mm}
{\small Метки: \textcolor{dzcol}{\textbf{ДЗ}} — задача из домашнего задания,
\textcolor{excol}{\textbf{пример}} — разобранный в книге пример (решён короче).
Формулы-«опоры» в начале каждого параграфа.}
\tableofcontents

\section*{§1. Первое начало термодинамики}
\addcontentsline{toc}{section}{§1. Первое начало термодинамики}

\begin{basics}
\sub{Уравнение состояния идеального газа (Клапейрон–Менделеев)}
\[pV=\nu RT=\frac m\mu RT=NkT,\qquad p=nkT.\]
$p$ — давление (Па), $V$ — объём (м$^3$), $T$ — абсолютная температура (К);
$\nu=m/\mu$ — число молей, $m$ — масса газа, $\mu$ — молярная масса (кг/моль);
$N$ — число молекул, $n=N/V$ — концентрация (м$^{-3}$);
$R=8{,}31$ Дж/(моль\,К), $k=R/N_A=1{,}38\cdot10^{-23}$ Дж/К, $N_A=6{,}02\cdot10^{23}$ моль$^{-1}$.
Смесь: $p=\sum p_i$ (закон Дальтона), $p_i$ — давление компоненты, если бы она была одна.

\sub{Внутренняя энергия идеального газа}
\[U=\frac i2\nu RT=\nu C_VT,\qquad C_V=\frac i2R.\]
$i$ — число степеней свободы молекулы: 3 (одноатомный: He, Ar), 5 (двухатомный: H$_2$, N$_2$, O$_2$, воздух), 6 (многоатомный).
$U$ зависит \emph{только от $T$} — поэтому $\Delta U=\nu C_V\Delta T$ в \emph{любом} процессе, не только изохорном.

\sub{Работа газа}
\[\delta A=p\,\dd V,\qquad A=\int_{V_1}^{V_2}p(V)\,\dd V\]
— площадь под графиком на $p$--$V$. При расширении $A>0$, при сжатии $A<0$. Зависит от пути.

\sub{Первое начало термодинамики}
\[\delta Q=\dd U+\delta A,\qquad Q=\Delta U+A.\]
$Q$ — теплота, \emph{полученная} газом ($Q<0$ — газ отдал тепло). Это закон сохранения энергии.
$\dd U$ — полный дифференциал (функция состояния), $\delta Q$ и $\delta A$ — нет (зависят от пути).
За цикл $\Delta U=0\Rightarrow A=\sum Q$.

\sub{Теплоёмкость}
$\tilde C=\delta Q/\dd T$ — теплоёмкость тела; молярная $C=\tilde C/\nu$, удельная $c=\tilde C/m$.
Зависит от процесса! Для моля идеального газа из первого начала:
\[C=\frac{\delta Q}{\dd T}=C_V+p\frac{\dd V}{\dd T}\qquad(\text{производная — вдоль процесса}).\]
Изобара: $p\frac{\dd V}{\dd T}=R\Rightarrow C_p=C_V+R$ (уравнение Майера).
Показатель адиабаты $\gamma=\dfrac{C_p}{C_V}=\dfrac{i+2}{i}$: $\frac53$, $\frac75$, $\frac43$; отсюда $C_V=\dfrac{R}{\gamma-1}$, $C_p=\dfrac{\gamma R}{\gamma-1}$.

\sub{Изопроцессы (на $\nu$ молей)}
\begin{tabular}{@{}llll@{}}
процесс & $A$ & $Q$ & $C$\\\hline
изохора $V$=const & $0$ & $\nu C_V\Delta T$ & $C_V$\\
изобара $p$=const & $p\Delta V=\nu R\Delta T$ & $\nu C_p\Delta T$ & $C_p$\\
изотерма $T$=const & $\nu RT\ln\frac{V_2}{V_1}$ & $Q=A$ ($\Delta U=0$) & $\infty$\\
адиабата $Q=0$ & $-\Delta U=\nu C_V(T_1-T_2)$ & $0$ & $0$
\end{tabular}

\sub{Адиабата (обратимая)}
\[pV^\gamma=\text{const},\qquad TV^{\gamma-1}=\text{const},\qquad p^{1-\gamma}T^\gamma=\text{const},\]
\[A=\frac{p_1V_1-p_2V_2}{\gamma-1}=\nu C_V(T_1-T_2).\]
Получается из $C=0$: $C_V\dd T+p\,\dd V=0$. Адиабата на $p$--$V$ круче изотермы.

\sub{Политропа ($C=\text{const}$, обобщает все изопроцессы)}
\[pV^n=\text{const},\qquad n=\frac{C-C_p}{C-C_V},\qquad C=C_V\frac{n-\gamma}{n-1},\]
\[A=\frac{p_1V_1-p_2V_2}{n-1},\qquad Q=\nu C\Delta T.\]

\sub{Циклы}
$Q_1$ — тепло, полученное от нагревателя (сумма всех $Q_i>0$), $Q_2$ — отданное холодильнику ($Q_i<0$).
\[\eta=\frac{A}{Q_1}=1-\frac{|Q_2|}{Q_1},\qquad\text{Карно: }\eta=1-\frac{T_2}{T_1};\]
\[\text{холодильник: }\eta_x=\frac{|Q_2|}{A}\ (\text{может быть}>1).\]
Цикл по часовой стрелке на $p$--$V$ — тепловая машина ($A>0$), против — холодильная.
\end{basics}

%-------------------------------------------------------------
\begin{problem}{1.1}{exdz}{С одним молем идеального газа проводят процесс $p=p_0-\beta V^2$, где $p_0$ и $\beta>0$ — постоянные.
Найти максимально возможную температуру газа в этом процессе.}
\textbf{1. Выражаем $T$ через $V$.} Параметры $p,V,T$ связаны уравнением состояния $pV=RT$ и уравнением процесса, поэтому независим один — берём $V$:
\[T(V)=\frac{pV}{R}=\frac{(p_0-\beta V^2)V}{R}=\frac{p_0V-\beta V^3}{R}.\]

\textbf{2. Ищем максимум:} $\frac{\dd T}{\dd V}=\frac{p_0-3\beta V^2}{R}=0\Rightarrow V_m=\sqrt{\frac{p_0}{3\beta}}$.

\textbf{3. Подставляем:} $p(V_m)=p_0-\frac{p_0}3=\frac{2p_0}3$, $T_{\max}=\frac{p(V_m)V_m}{R}$.
\answer{$T_{\max}=\dfrac{2p_0}{3R}\sqrt{\dfrac{p_0}{3\beta}}$ (при $V=\sqrt{p_0/3\beta}$).}
\note{так же: 1.14 ($p=p_0e^{-\beta V}$, $T_{\max}=\frac{p_0}{\beta eR}$), 1.15 (минимум $p$).}
\end{problem}

%-------------------------------------------------------------
\begin{problem}{1.2}{ex}{Для идеального газа при $0^\circ$C найти термические коэффициенты: расширения $\alpha=\frac1{V_0}\big(\frac{\partial V}{\partial T}\big)_p$,
давления $\lambda=\frac1{p_0}\big(\frac{\partial p}{\partial T}\big)_V$ и сжатия $\beta=-\frac1{V_0}\big(\frac{\partial V}{\partial p}\big)_T$. Индекс 0 — при $T_0=273$\,К.}
Дифференцируем $pV=\nu RT$, держа постоянным указанный параметр:
\begin{itemize}
\item $p=$\,const: $p\,\dd V=\nu R\,\dd T\Rightarrow\big(\frac{\partial V}{\partial T}\big)_p=\frac{\nu R}{p}=\frac{V}{T}$, т.е.\ $\alpha=\frac1{T_0}$;
\item $V=$\,const: $\big(\frac{\partial p}{\partial T}\big)_V=\frac{\nu R}{V}=\frac pT$, т.е.\ $\lambda=\frac1{T_0}$;
\item $T=$\,const: $V=\frac{\nu RT}p\Rightarrow\big(\frac{\partial V}{\partial p}\big)_T=-\frac{V}{p}$, т.е.\ $\beta=\frac1{p_0}$.
\end{itemize}
\answer{$\alpha=\lambda=1/T_0$,\; $\beta=1/p_0$; выполняется $\alpha=p_0\beta\lambda$.}
\end{problem}

%-------------------------------------------------------------
\begin{problem}{1.3}{exdz}{Определить массу воздуха между двумя оконными рамами при атмосферном давлении $p_0=10^5$\,Па, если температура между рамами
меняется по линейному закону от $T_1=270$\,К до $T_2=300$\,К. Площадь окна $S=3$\,м$^2$, расстояние между рамами $l=0{,}2$\,м, $\mu=29$\,г/моль.}
\textbf{Идея:} газ неоднороден, но тонкий слой $\dd x$ можно считать однородным.

\textbf{1.} Ось $x$ перпендикулярно рамам: $T(x)=T_1+\frac{T_2-T_1}{l}x$. Для слоя толщиной $\dd x$ (объём $S\,\dd x$) уравнение состояния:
$p_0S\,\dd x=\frac{\dd m}{\mu}RT(x)\Rightarrow\dd m=\frac{\mu p_0S}{R}\cdot\frac{\dd x}{T(x)}$.

\textbf{2. Интегрируем} от 0 до $l$ (интеграл вида $\int\frac{\dd x}{a+bx}=\frac1b\ln(a+bx)$):
\[m=\frac{\mu p_0S}{R}\int_0^l\frac{\dd x}{T_1+(T_2-T_1)x/l}=\frac{\mu p_0Sl}{R(T_2-T_1)}\ln\frac{T_2}{T_1}.\]

\textbf{3. Числа:} $\frac{0{,}029\cdot10^5\cdot3\cdot0{,}2}{8{,}31\cdot30}=6{,}98$\,кг, $\ln\frac{300}{270}=0{,}105$.
\answer{$m=\dfrac{\mu p_0Sl}{R(T_2-T_1)}\ln\dfrac{T_2}{T_1}\approx0{,}735$ кг.}
\note{при $T_2\to T_1$ получаем обычное $m=\frac{\mu pV}{RT}$ — проверка.}
\end{problem}

%-------------------------------------------------------------
\begin{problem}{1.4}{ex}{Моль идеального газа ($i=5$) три раза обратимо переводят из состояния 1 $(p_1,V_1)$ в состояние 2 $(2p_1,2V_1)$ путями 1--3--2, 1--4--2 и 1--2
(рис.~1.5: 3 — $(2p_1,V_1)$, 4 — $(p_1,2V_1)$, 1--2 — отрезок прямой). Найти $Q_{132}$, $Q_{142}$, $Q_{12}$ и молярную теплоёмкость $C_{12}$. Выразить через $T_1$.}
\textbf{1. Температуры} из $pV=RT$: $T_3=T_4=2T_1$, $T_2=4T_1$. Для $i=5$: $C_V=\frac52R$, $C_p=\frac72R$.

\textbf{2. Путь 1--3--2:} изохора 1--3, $Q=C_V(2T_1-T_1)=2{,}5RT_1$; изобара 3--2, $Q=C_p(4T_1-2T_1)=7RT_1$. Сумма $9{,}5RT_1$.

\textbf{3. Путь 1--4--2:} изобара 1--4, $Q=C_p\cdot T_1=3{,}5RT_1$; изохора 4--2, $Q=C_V\cdot2T_1=5RT_1$. Сумма $8{,}5RT_1$.

\textbf{4. Путь 1--2:} работа — площадь трапеции под отрезком: $A=\frac{p_1+2p_1}2V_1=1{,}5RT_1$; $\Delta U=C_V(4T_1-T_1)=7{,}5RT_1$; $Q=\Delta U+A=9RT_1$.

\textbf{5. Теплоёмкость на 1--2.} Прямая через начало координат: $p=\alpha V$. Тогда $RT=pV=\alpha V^2$, дифференцируем: $R\,\dd T=2\alpha V\dd V=2p\,\dd V$,
т.е.\ $p\frac{\dd V}{\dd T}=\frac R2$ и $C=C_V+p\frac{\dd V}{\dd T}=\frac52R+\frac12R$.
\answer{$Q_{132}=9{,}5RT_1$, $Q_{142}=8{,}5RT_1$, $Q_{12}=9RT_1$, $C_{12}=3R$.}
\end{problem}

%-------------------------------------------------------------
\begin{problem}{1.5}{exdz}{Политропным называют обратимый процесс при постоянной теплоёмкости $C$. Найти уравнение политропы для моля идеального газа с известной $C_V$.}
\textbf{1. Теплоёмкость из первого начала:} $C\,\dd T=C_V\dd T+p\,\dd V$, $p=\frac{RT}V$:
\[(C-C_V)\,\dd T=\frac{RT}{V}\dd V\ \Rightarrow\ (C-C_V)\frac{\dd T}{T}=R\frac{\dd V}{V}.\]

\textbf{2. Интегрируем} ($C=$\,const): $(C-C_V)\ln T=R\ln V+\text{const}\Rightarrow T^{C_V-C}V^R=\text{const}$.

\textbf{3. Переходим к $p,V$:} возводим в степень $\frac1{C_V-C}$: $TV^{R/(C_V-C)}=\text{const}$; подставляем $T=\frac{pV}R$:
$pV^{n}=\text{const}$, где $n=1+\frac{R}{C_V-C}=\frac{C_V+R-C}{C_V-C}=\frac{C-C_p}{C-C_V}$.
\answer{$pV^n=\text{const}$ (и $TV^{n-1}=\text{const}$), $n=\dfrac{C-C_p}{C-C_V}$; обратно $C=C_V\dfrac{n-\gamma}{n-1}$.}
\note{$n=0$ — изобара ($C_p$), $n=1$ — изотерма ($C=\infty$), $n=\gamma$ — адиабата ($C=0$), $n=\infty$ — изохора ($C_V$); при $1<n<\gamma$ теплоёмкость $C<0$.}
\end{problem}

%-------------------------------------------------------------
\begin{problem}{1.6}{exdz}{Боковые стенки цилиндра, его крышка и невесомый поршень не проводят тепло, дно проводит. Над и под поршнем — по молю одного и того же
идеального газа с известной $C_V$. Газ в нижней части нагревают, поршень перемещается. Найти теплоёмкости $C_1$ (нижний газ) и $C_2$ (верхний) через их объёмы $V_1$, $V_2$.}
\textbf{1. Верхний газ} теплоизолирован, процесс медленный — адиабата: $C_2=0$ и $pV_2^\gamma=\text{const}$.
Дифференцируем логарифм: $\frac{\dd p}{p}=-\gamma\frac{\dd V_2}{V_2}$. Объём цилиндра постоянен: $\dd V_2=-\dd V_1$, значит $\frac{\dd p}{p}=\gamma\frac{\dd V_1}{V_2}$.

\textbf{2. Нижний газ.} Давления по обе стороны невесомого поршня равны. Уравнение состояния $pV_1=RT_1$, дифференцируем:
\[p\,\dd V_1+V_1\dd p=R\,\dd T_1\ \Rightarrow\ p\,\dd V_1\Big(1+\gamma\frac{V_1}{V_2}\Big)=R\,\dd T_1\ \Rightarrow\ p\frac{\dd V_1}{\dd T_1}=\frac{RV_2}{V_2+\gamma V_1}.\]

\textbf{3. Теплоёмкость} $C_1=C_V+p\frac{\dd V_1}{\dd T_1}$.
\answer{$C_1=C_V+R\dfrac{V_2}{V_2+\gamma V_1}$,\; $C_2=0$.}
\end{problem}

%-------------------------------------------------------------
\begin{problem}{1.7}{ex}{Найти разность теплоёмкостей $C_p-C_V$ для любого однородного изотропного вещества с известными уравнением состояния $p(V,T)$ и внутренней энергией $U(T,V)$.}
\textbf{1. Первое начало} с полным дифференциалом $U$:
$\delta Q=\dd U+p\,\dd V=\big(\frac{\partial U}{\partial T}\big)_V\dd T+\big[\big(\frac{\partial U}{\partial V}\big)_T+p\big]\dd V$.

\textbf{2.} Делим на $\dd T$. При $V=$\,const второе слагаемое пропадает: $C_V=\big(\frac{\partial U}{\partial T}\big)_V$.
При $p=$\,const: $C_p=C_V+\big[\big(\frac{\partial U}{\partial V}\big)_T+p\big]\big(\frac{\partial V}{\partial T}\big)_p$.
\answer{$C_p-C_V=\Big[\Big(\dfrac{\partial U}{\partial V}\Big)_T+p\Big]\Big(\dfrac{\partial V}{\partial T}\Big)_p$; для идеального газа ($\frac{\partial U}{\partial V}=0$, $p\frac{\partial V}{\partial T}=R$) это $R$.}
\end{problem}

%-------------------------------------------------------------
\begin{problem}{1.8}{exdz}{Обратимый цикл с молем идеального газа (рис.~1.9): 1--2 — изобара ($T_1\to T_2$), 2--3 — изохора ($T_2\to T_3$), 3--1 — адиабата.
Найти работу и теплоту на каждом этапе и КПД цикла как функцию $T_1$, $T_2$, $T_3$. $\gamma$ известна.}
\textbf{1. По участкам} (первое начало $Q=\Delta U+A$, $\Delta U=C_V\Delta T$):
\begin{itemize}
\item 1--2, изобара: $A_{12}=p_1(V_2-V_1)=R(T_2-T_1)$, $Q_{12}=C_p(T_2-T_1)>0$ — тепло получено;
\item 2--3, изохора: $A_{23}=0$, $Q_{23}=C_V(T_3-T_2)<0$ — тепло отдано;
\item 3--1, адиабата: $Q_{31}=0$, $A_{31}=-\Delta U=C_V(T_3-T_1)$.
\end{itemize}

\textbf{2. КПД:} $A=Q_{12}+Q_{23}$, тепло получено только на 1--2:
\[\eta=\frac{A}{Q_{12}}=1-\frac{C_V(T_2-T_3)}{C_p(T_2-T_1)}.\]
\answer{$\eta=1-\dfrac{T_2-T_3}{\gamma(T_2-T_1)}$.}
\end{problem}

%-------------------------------------------------------------
\begin{problem}{1.9}{ex}{Найти КПД цикла Карно (две изотермы $T_1>T_2$ и две адиабаты) с идеальным газом.}
\textbf{1. Теплоты на изотермах} ($\Delta U=0$, $Q=A$): $Q_{12}=\nu RT_1\ln\frac{V_2}{V_1}>0$, $Q_{34}=\nu RT_2\ln\frac{V_4}{V_3}<0$.

\textbf{2. Связь объёмов} — из адиабат 2--3 и 4--1 ($TV^{\gamma-1}=$\,const): $T_1V_2^{\gamma-1}=T_2V_3^{\gamma-1}$, $T_1V_1^{\gamma-1}=T_2V_4^{\gamma-1}$.
Делим одно на другое: $\frac{V_2}{V_1}=\frac{V_3}{V_4}$.

\textbf{3. КПД:} $\eta=1+\frac{Q_{34}}{Q_{12}}=1-\frac{T_2\ln(V_3/V_4)}{T_1\ln(V_2/V_1)}$, логарифмы сокращаются.
\answer{$\eta=1-\dfrac{T_2}{T_1}$.}
\end{problem}

%-------------------------------------------------------------
\begin{problem}{1.10}{ex}{Машина Карно с КПД 40\% используется как холодильная при тех же резервуарах. Какое количество теплоты $Q_2$ она переносит от холодильника
к нагревателю за цикл, если за цикл совершается работа $A=10$\,кДж? Оценить её эффективность.}
\textbf{1.} При обращении цикла все теплоты меняют знак, но по модулю те же. Для тепловой машины: $\eta=\frac A{Q_1}$, $A=Q_1-|Q_2|$.

\textbf{2.} Исключаем $Q_1=\frac A\eta$: $|Q_2|=Q_1-A=A\big(\frac1\eta-1\big)=10\cdot(2{,}5-1)=15$\,кДж.

\textbf{3. Холодильный коэффициент} $\eta_x=\frac{|Q_2|}{A}$.
\answer{$|Q_2|=15$ кДж,\; $\eta_x=1{,}5$.}
\end{problem}

%-------------------------------------------------------------
\begin{problem}{1.25}{dz}{Найти уравнение процесса для идеального газа, при котором молярная теплоёмкость зависит от температуры как $C=\alpha T$, $\alpha$ — постоянная. $C_V$ известна.}
\textbf{1. Теплоёмкость в процессе} (как в 1.5): $C=C_V+p\frac{\dd V}{\dd T}=C_V+\frac{RT}{V}\frac{\dd V}{\dd T}$. Приравниваем $\alpha T$:
\[\alpha T-C_V=\frac{RT}{V}\frac{\dd V}{\dd T}\ \Rightarrow\ R\frac{\dd V}{V}=\alpha\,\dd T-C_V\frac{\dd T}{T}.\]

\textbf{2. Интегрируем:} $R\ln V=\alpha T-C_V\ln T+\text{const}$. Делим на $R$ и потенцируем: $V\cdot T^{C_V/R}\cdot e^{-\alpha T/R}=\text{const}$.
\answer{$VT^{C_V/R}\exp(-\alpha T/R)=\text{const}$.}
\end{problem}

%-------------------------------------------------------------
\begin{problem}{1.28}{dz}{Моль идеального газа нагревают в цилиндре под поршнем, удерживаемым пружиной (закон Гука). Стенки и поршень теплоизолированы, дно проводит тепло.
Начальный объём $V_0$ (пружина не деформирована) подобран так, что $p_0S^2=kV_0$, где $p_0$ — наружное атмосферное давление, $S$ — площадь поршня,
$k$ — жёсткость пружины. Найти теплоёмкость газа. Массой поршня пренебречь.}
\textbf{1. Уравнение процесса.} Поршень сместился на $x=\frac{V-V_0}S$, пружина давит с силой $kx$. Равновесие поршня:
\[pS=p_0S+kx\ \Rightarrow\ p=p_0+\frac{k}{S^2}(V-V_0)=p_0+\frac{p_0}{V_0}(V-V_0)=\frac{p_0}{V_0}V\]
(использовали $\frac k{S^2}=\frac{p_0}{V_0}$). Процесс $p=\alpha V$ — прямая через начало координат.

\textbf{2. Теплоёмкость.} Как в 1.4: $RT=\alpha V^2\Rightarrow R\,\dd T=2p\,\dd V\Rightarrow p\frac{\dd V}{\dd T}=\frac R2$.
\answer{$C=C_V+\dfrac R2$.}
\end{problem}

%-------------------------------------------------------------
\begin{problem}{1.33}{dz}{Обратимый цикл с молем идеального газа: две изотермы 1--2 ($T_1$) и 3--4 ($T_2$) и две политропы 2--3, 4--1 с теплоёмкостью $C_0$ (рис.~1.15).
Найти КПД, если при изотермическом расширении объём увеличивается вдвое.}
\textbf{1. Отношение объёмов на изотермах.} На политропе $T^{C_V-C_0}V^R=$\,const (см.\ 1.5), т.е.\ $V\propto T^{(C_0-C_V)/R}$.
Обе политропы идут между $T_1$ и $T_2$, поэтому $\frac{V_3}{V_2}=\frac{V_4}{V_1}$, откуда $\frac{V_3}{V_4}=\frac{V_2}{V_1}=2$.

\textbf{2. Теплоты по участкам:}
\begin{itemize}
\item 1--2 (расширение при $T_1$): $Q=RT_1\ln2>0$;
\item 2--3 (политропа $T_1\to T_2$): $Q=C_0(T_2-T_1)<0$;
\item 3--4 (сжатие при $T_2$): $Q=-RT_2\ln2<0$;
\item 4--1 (политропа $T_2\to T_1$): $Q=C_0(T_1-T_2)>0$ (при $C_0>0$).
\end{itemize}

\textbf{3. КПД.} Работа $A=\sum Q=R(T_1-T_2)\ln2$ (теплоты на политропах сокращаются). Получено тепла $Q_+=RT_1\ln2+C_0(T_1-T_2)$.
\answer{$\eta=\dfrac{R(T_1-T_2)\ln2}{C_0(T_1-T_2)+RT_1\ln2}$.}
\note{1.32 (вместо политроп — изобары) — то же с $C_0\to C_p$.}
\end{problem}

%-------------------------------------------------------------
\begin{problem}{1.34}{dz}{Определить, каким путём теоретически эффективнее повысить КПД машины Карно: увеличить температуру нагревателя $T_1$ на $\Delta T$ при
фиксированной $T_2$ или понизить температуру холодильника $T_2$ на такую же $\Delta T$ при фиксированной $T_1$.}
\textbf{1. Два варианта:} $\eta_a=1-\dfrac{T_2}{T_1+\Delta T}$ (греем нагреватель), $\eta_b=1-\dfrac{T_2-\Delta T}{T_1}$ (охлаждаем холодильник).

\textbf{2. Сравниваем} дроби: $\eta_b>\eta_a\iff\frac{T_2-\Delta T}{T_1}<\frac{T_2}{T_1+\Delta T}\iff(T_2-\Delta T)(T_1+\Delta T)<T_1T_2$.
Раскрываем скобки: $T_2\Delta T-T_1\Delta T-\Delta T^2<0$, т.е.\ $-\Delta T(T_1-T_2+\Delta T)<0$ — верно, так как $T_1>T_2$.
\answer{эффективнее понижать температуру холодильника.}
\end{problem}

%-------------------------------------------------------------
\begin{problem}{1.35}{dz}{Моль одноатомного идеального газа ($\gamma=5/3$) совершает цикл Карно между резервуарами $t_1=127^\circ$C и $t_2=27^\circ$C.
Наименьший объём газа в цикле $V_1=5$\,л, наибольший $V_2=20$\,л. Найти работу за цикл $A$, теплоту $Q_1$ от нагревателя и теплоту $Q_2$, отданную холодильнику.}
\textbf{1. Где какие объёмы.} Наименьший объём — начало изотермического расширения при $T_1=400$\,К. Наибольший — конец расширения по адиабате, т.е.\ уже при $T_2=300$\,К.
Объём $V_a$ в конце изотермы $T_1$ находим по адиабате $V_a\to V_2$:
\[T_1V_a^{\gamma-1}=T_2V_2^{\gamma-1}\ \Rightarrow\ V_a=V_2\Big(\frac{T_2}{T_1}\Big)^{1/(\gamma-1)}=20\cdot0{,}75^{3/2}\approx13\text{ л}.\]

\textbf{2. Тепло от нагревателя} (изотерма $T_1$, $V_1\to V_a$):
$Q_1=RT_1\ln\frac{V_a}{V_1}=8{,}31\cdot400\cdot\ln2{,}6\approx3{,}17$\,кДж.

\textbf{3. Карно:} $\frac{|Q_2|}{Q_1}=\frac{T_2}{T_1}\Rightarrow Q_2=-\frac{300}{400}Q_1=-2{,}38$\,кДж; $A=Q_1+Q_2$.
\answer{$Q_1\approx3{,}17$ кДж, $Q_2\approx-2{,}38$ кДж, $A\approx0{,}79$ кДж.}
\end{problem}

%-------------------------------------------------------------
\begin{problem}{1.36}{dz}{Определить КПД цикла 1--2--3--4--1, в котором одноатомный идеальный газ последовательно проходит состояния $1(2p,V)$, $2(2p,8V)$, $3(p,4V)$, $4(p,2V)$.
Все участки цикла — политропы.}
\textbf{1. Температуры} в единицах $\frac{pV}{\nu R}$ (произведение координат): $T_1=2$, $T_2=16$, $T_3=4$, $T_4=2$.

\textbf{2. Тип каждого участка} (из $p_aV_a^n=p_bV_b^n$) и теплота $Q=\nu C\Delta T$, для одноатомного $C_V=\frac32R$, $C_p=\frac52R$:
\begin{itemize}
\item 1--2: $p=2p$ — изобара; $Q=\frac52\cdot(16-2)=35$;
\item 2--3: $\frac{p}{V}$ одинаково ($\frac{2p}{8V}=\frac{p}{4V}$) — процесс $p=\alpha V$, $C=C_V+\frac R2=2R$; $Q=2\cdot(4-16)=-24$;
\item 3--4: $p$ — изобара; $Q=\frac52\cdot(2-4)=-5$;
\item 4--1: $pV=2pV$ одинаково — изотерма, сжатие $2V\to V$; $Q=RT\ln\frac{V}{2V}=2\ln\frac12=-2\ln2$.
\end{itemize}
(все $Q$ в единицах $pV$.)

\textbf{3. КПД.} $A=\sum Q=35-24-5-2\ln2=6-2\ln2$; тепло получено только на 1--2: $Q_+=35$.
\answer{$\eta=\dfrac{6-2\ln2}{35}\approx13\%$.}
\end{problem}

%-------------------------------------------------------------
\begin{problem}{1.37}{dz}{Найти КПД обратимого цикла Отто из адиабат 1--2, 3--4 и изохор 2--3, 4--1 (рис.~1.16), рабочее тело — идеальный газ. Выразить через $T_1$ и $T_2$.}
\textbf{1. Теплоты.} На адиабатах $Q=0$. Изохора 4--1 — нагрев, тепло получено: $Q_1=C_V(T_1-T_4)$. Изохора 2--3 — охлаждение: $|Q_2|=C_V(T_2-T_3)$.

\textbf{2. Связь температур.} Обе адиабаты идут между одними и теми же объёмами $V_{\min}$ и $V_{\max}$ ($TV^{\gamma-1}=$\,const):
$\frac{T_2}{T_1}=\frac{T_3}{T_4}=\big(\frac{V_{\min}}{V_{\max}}\big)^{\gamma-1}$. Значит $T_2=kT_1$, $T_3=kT_4$ и $\frac{T_2-T_3}{T_1-T_4}=k=\frac{T_2}{T_1}$.

\textbf{3.} $\eta=1-\frac{|Q_2|}{Q_1}=1-\frac{T_2-T_3}{T_1-T_4}$.
\answer{$\eta=1-\dfrac{T_2}{T_1}$.}
\end{problem}

\section*{§2. Второе начало. Энтропия}
\addcontentsline{toc}{section}{§2. Второе начало. Энтропия}

\begin{basics}
\sub{Энтропия}
\[\dd S=\frac{\delta Q}{T}\quad\text{(только для обратимого процесса)},\qquad S_2-S_1=\int_1^2\frac{\delta Q}{T}.\]
$S$ — функция состояния (Дж/К): $\Delta S$ зависит только от начального и конечного состояний, не от пути.
$1/T$ — интегрирующий множитель: $\delta Q$ — не полный дифференциал, а $\delta Q/T$ — полный.

\sub{Второе начало}
\begin{itemize}
\item Неравенство Клаузиуса: $\displaystyle\oint\frac{\delta Q}{T}\le0$; «$=$» для обратимого цикла, «$<$» для необратимого.
\item В изолированной (теплоизолированной) системе $\Delta S\ge0$: при необратимых процессах (смешение, теплообмен при
разных $T$, трение, расширение в пустоту, резкий скачок давления) энтропия \emph{растёт}, при обратимых — постоянна.
\item КПД любой машины между $T_1$ и $T_2$ не больше, чем у Карно: $\eta\le1-T_2/T_1$.
\end{itemize}

\sub{Как считать $\Delta S$ в необратимом процессе}
Формула $\dd S=\delta Q/T$ применима только к обратимому пути. Поэтому:
(1) найти конечное равновесное состояние (из закона сохранения энергии и условий равновесия);
(2) соединить начало и конец \emph{любым удобным обратимым} путём и посчитать $\Delta S$ по нему.
Для идеального газа это уже сделано в формуле ниже.

\sub{Энтропия идеального газа}
Из $T\dd S=\nu C_V\dd T+p\,\dd V$, $p=\nu RT/V$:
\begin{align*}
\Delta S&=\nu\Big[C_V\ln\frac{T_2}{T_1}+R\ln\frac{V_2}{V_1}\Big]
=\nu\Big[C_p\ln\frac{T_2}{T_1}-R\ln\frac{p_2}{p_1}\Big]\\
&=\nu\Big[C_V\ln\frac{p_2}{p_1}+C_p\ln\frac{V_2}{V_1}\Big].
\end{align*}
Частные случаи: изотерма $\nu R\ln\frac{V_2}{V_1}$, изохора $\nu C_V\ln\frac{T_2}{T_1}$,
изобара $\nu C_p\ln\frac{T_2}{T_1}$, обратимая адиабата $0$, политропа $\nu C\ln\frac{T_2}{T_1}$.
Энтропия аддитивна: у смеси/системы $\Delta S=\sum_i\Delta S_i$, каждая компонента считается так, будто она одна
занимает весь свой объём.

\sub{Конденсированные тела и фазовые переходы}
Нагрев тела с теплоёмкостью $C$ ($=mc$): $\Delta S=\int\frac{C\,\dd T}{T}=C\ln\frac{T_2}{T_1}$.
\\Плавление/кипение при постоянной $T$: $\Delta S=\frac{Q}{T}=\frac{\lambda m}{T}$, $\lambda$ — удельная теплота перехода.

\sub{Диаграмма $T$--$S$}
$Q=\int T\,\dd S$ — площадь под кривой; $A$ за цикл — площадь цикла; изотерма — горизонталь, обратимая адиабата — вертикаль.

\sub{Полезное тождество}
$\Big(\dfrac{\partial U}{\partial V}\Big)_T=T\Big(\dfrac{\partial p}{\partial T}\Big)_V-p$ (из $\dd U=T\dd S-p\,\dd V$); для идеального газа даёт $0$.
\end{basics}

%-------------------------------------------------------------
\begin{problem}{2.1}{ex}{Азот, занимающий объём $V_1=1$\,м$^3$ при давлении $p=10^5$\,Па и температуре $T=300$\,К, расширяется изотермически
до объёма $V_2=2$\,м$^3$. Вычислить изменение энтропии газа $\Delta S$.}
\textbf{1.} Изотерма: $\Delta S=\nu R\ln\frac{V_2}{V_1}$ (формула энтропии ид.\ газа при $T_2=T_1$).

\textbf{2.} Число молей из уравнения состояния: $\nu R=\frac{pV_1}{T}$. Тогда
$\Delta S=\frac{pV_1}{T}\ln\frac{V_2}{V_1}=\frac{10^5\cdot1}{300}\cdot\ln2=333\cdot0{,}693$.
\answer{$\Delta S\approx231$ Дж/К.}
\end{problem}

%-------------------------------------------------------------
\begin{problem}{2.2}{ex}{Найти изменение энтропии льда массой $m=0{,}03$\,кг при превращении его в пар, если начальная температура льда
$-40^\circ$C, а температура пара $100^\circ$C. Удельные теплоёмкости льда и воды $c_1=2{,}1$ и $c_2=4{,}2$\,кДж/(кг\,К);
теплота плавления $\lambda_1=330$\,кДж/кг, парообразования $\lambda_2=2{,}25\cdot10^3$\,кДж/кг. Давление атмосферное.}
\textbf{Идея:} разбить процесс на этапы и сложить изменения энтропии (энтропия аддитивна).
\begin{enumerate}
\item Нагрев льда $233\to273$\,К: $\Delta S_1=mc_1\ln\frac{273}{233}=0{,}03\cdot2100\cdot0{,}158=10$\,Дж/К.
\item Плавление при $273$\,К: $\Delta S_2=\frac{m\lambda_1}{273}=\frac{0{,}03\cdot3{,}3\cdot10^5}{273}=36{,}2$\,Дж/К.
\item Нагрев воды $273\to373$\,К: $\Delta S_3=mc_2\ln\frac{373}{273}=0{,}03\cdot4200\cdot0{,}312=39{,}3$\,Дж/К.
\item Кипение при $373$\,К: $\Delta S_4=\frac{m\lambda_2}{373}=\frac{0{,}03\cdot2{,}25\cdot10^6}{373}=181$\,Дж/К.
\end{enumerate}
\answer{$\Delta S=\Delta S_1+\Delta S_2+\Delta S_3+\Delta S_4\approx266{,}5$ Дж/К.}
\end{problem}

%-------------------------------------------------------------
\begin{problem}{2.3}{exdz}{Определить работу, совершаемую произвольным телом за обратимый цикл из изотермы 1--2, политропы 2--3 и адиабаты 3--1.
Известны теплоёмкость $C$ тела в политропном процессе и температуры $T_1$ (на изотерме) и $T_3$.}
\textbf{Идея:} уравнение состояния тела неизвестно, поэтому теплоту на изотерме находим из второго начала.

\textbf{1. Равенство Клаузиуса} для обратимого цикла $\oint\frac{\delta Q}{T}=0$, по участкам:
\[\underbrace{\frac{Q_{12}}{T_1}}_{\text{изотерма}}+\underbrace{\int_{T_1}^{T_3}\frac{C\,\dd T}{T}}_{\text{политропа}}+\underbrace{0}_{\text{адиабата}}=0
\ \Rightarrow\ Q_{12}=-CT_1\ln\frac{T_3}{T_1}=CT_1\ln\frac{T_1}{T_3}.\]

\textbf{2. Теплота на политропе:} $Q_{23}=C(T_3-T_1)$; на адиабате $Q_{31}=0$.

\textbf{3. Работа за цикл} $A=\sum Q$ (за цикл $\Delta U=0$).
\answer{$A=C\Big[T_1\ln\dfrac{T_1}{T_3}+T_3-T_1\Big]$.}
\end{problem}

%-------------------------------------------------------------
\begin{problem}{2.4}{ex}{Показать, что для любого вещества политропа может пересекать изотерму не более чем в одной точке.}
\textbf{От противного.} Пусть политропа и изотерма $T_0$ пересекаются в точках $A$ и $B$. Составим цикл: $A\to B$ по политропе, $B\to A$ по изотерме.
\begin{enumerate}
\item На политропе ($C=$\,const): $Q_{\text{пол}}=C(T_B-T_A)=0$ и $\int\frac{C\,\dd T}{T}=C\ln\frac{T_B}{T_A}=0$, так как $T_A=T_B=T_0$.
\item Клаузиус для обратимого цикла: $0+\frac{Q_{\text{из}}}{T_0}=0\Rightarrow Q_{\text{из}}=0$.
\item Тогда работа за цикл $A=Q_{\text{пол}}+Q_{\text{из}}=0$. Но кривые разные, цикл охватывает ненулевую площадь на $p$--$V$, т.е.\ $A\ne0$. Противоречие.
\end{enumerate}
\end{problem}

%-------------------------------------------------------------
\begin{problem}{2.5}{ex}{Электромагнитное излучение в полости объёмом $V$ находится в равновесии со стенками при температуре $T$.
Его энергия $U=u(T)V$, давление $p=u/3$. Найти связь $u$ с $T$ (закон Стефана–Больцмана).}
\textbf{1. Тождество} из $\dd U=T\dd S-p\,\dd V$ (следствие второго начала):
$\big(\frac{\partial U}{\partial V}\big)_T=T\big(\frac{\partial p}{\partial T}\big)_V-p$.

\textbf{2. Подставляем:} $\big(\frac{\partial U}{\partial V}\big)_T=u$, $\big(\frac{\partial p}{\partial T}\big)_V=\frac{u'}3$:
\quad $u=\frac T3u'-\frac u3\Rightarrow4u=T\frac{\dd u}{\dd T}$.

\textbf{3. Разделяем переменные:} $\frac{\dd u}{u}=4\frac{\dd T}{T}\Rightarrow\ln u=4\ln T+\text{const}$.
\answer{$u=\sigma T^4$.}
\end{problem}

%-------------------------------------------------------------
\begin{problem}{2.6}{exdz}{В качестве переменных состояния выберем $T$ и $S$. Построить цикл Карно на диаграмме $T$--$S$ и найти с её помощью КПД.}
\textbf{1. Вид цикла.} Изотермы $T_1$ и $T_2$ — горизонтальные прямые. Обратимые адиабаты: $\delta Q=0\Rightarrow\dd S=0$ — вертикальные прямые $S_1$ и $S_2$.
Цикл — прямоугольник.

\textbf{2. Теплоты.} На $T$--$S$ теплота $Q=\int T\,\dd S$ — площадь под кривой:
на верхней изотерме получено $Q_1=T_1(S_2-S_1)$, на нижней отдано $|Q_2|=T_2(S_2-S_1)$.

\textbf{3. КПД:} $\eta=\frac{Q_1-|Q_2|}{Q_1}$, $(S_2-S_1)$ сокращается.
\answer{$\eta=\dfrac{T_1-T_2}{T_1}$ — для любого рабочего тела.}
\end{problem}

%-------------------------------------------------------------
\begin{problem}{2.7}{ex}{Система квазистатически переводится из состояния 1 в состояние 2 двумя способами: (а) адиабатически охлаждается до $T_0$,
изотермически получает тепло, затем адиабатически переходит в 2; (б) по пути, на каждом участке которого система получает тепло, а её температура выше $T_0$.
В каком случае система получает больше тепла?}
\textbf{1.} Изменение энтропии $S_2-S_1$ одинаково (функция состояния), и $S_2>S_1$, так как в (б) всё время $\delta Q>0$.

\textbf{2. Способ (а):} тепло получено только на изотерме: $Q_a=T_0(S_2-S_1)$.

\textbf{3. Способ (б):} $Q_b=\int T\,\dd S$, и везде $T>T_0$, $\dd S>0$, поэтому $Q_b>T_0\int\dd S=T_0(S_2-S_1)$.
\answer{$Q_b>Q_a$.}
\end{problem}

%-------------------------------------------------------------
\begin{problem}{2.8}{ex}{Теплоизолированный сосуд разделён невесомым поршнем пополам. С одной стороны $\nu$ молей идеального газа при $T_0$, $p_0$,
с другой — вакуум. Поршень отпускают, газ заполняет весь сосуд. Затем объём медленно возвращают к исходному. Найти изменения внутренней энергии и энтропии.}
\textbf{1. Расширение в пустоту.} Поршень невесомый — газ работы не совершает ($A=0$), сосуд теплоизолирован ($Q=0$), значит $\Delta U=0$ и $T=T_0$.
Процесс необратимый, но $\Delta S$ считаем по изотерме между теми же состояниями: $\Delta S_1=\nu R\ln2$.

\textbf{2. Медленное сжатие} в теплоизолированном сосуде — обратимая адиабата: $\Delta S_2=0$;
$TV^{\gamma-1}=\text{const}\Rightarrow T=T_0\cdot2^{\gamma-1}$.

\textbf{3. Итог:} $\Delta U=\nu C_V(T-T_0)$, $\Delta S=\Delta S_1+\Delta S_2$.
\answer{$\Delta U=\nu C_VT_0\big(2^{\gamma-1}-1\big)$,\; $\Delta S=\nu R\ln2$.}
\end{problem}

%-------------------------------------------------------------
\begin{problem}{2.9}{exdz}{Приводимые в тепловой контакт два тела имеют одинаковые теплоёмкости $C_p=\text{const}$ и разные температуры $T_1$ и $T_2$.
Найти приращение энтропии в результате установления теплового равновесия.}
\textbf{1. Конечная температура} — из сохранения энергии: $C_p(T-T_1)+C_p(T-T_2)=0\Rightarrow T=\frac{T_1+T_2}2$.

\textbf{2. Энтропия каждого тела} при нагреве/охлаждении: $\Delta S_i=C_p\ln\frac{T}{T_i}$. Складываем:
$\Delta S=C_p\ln\frac{T^2}{T_1T_2}$.
\answer{$\Delta S=C_p\ln\dfrac{(T_1+T_2)^2}{4T_1T_2}$.}
\note{$\Delta S>0$, т.к.\ $(T_1+T_2)^2\ge4T_1T_2$ (среднее арифметическое $\ge$ среднего геометрического) — процесс необратим.}
\end{problem}

%-------------------------------------------------------------
\begin{problem}{2.10}{exdz}{В двух одинаковых сосудах находятся различные идеальные газы: масса первого $M_1$ (молярная масса $\mu_1$), второго — $M_2$ ($\mu_2$).
Температуры одинаковы. Сосуды соединяют, начинается диффузия. Найти суммарное изменение энтропии.}
\textbf{1. Конечное состояние.} Газы идеальные — не взаимодействуют; энергия сохраняется, температура не меняется.
Каждый газ просто расширяется на весь объём: $V\to2V$.

\textbf{2. Энтропия каждого газа} — как при изотермическом расширении: $\Delta S_i=\nu_iR\ln2=\frac{M_i}{\mu_i}R\ln2$.
\answer{$\Delta S=\Big(\dfrac{M_1}{\mu_1}+\dfrac{M_2}{\mu_2}\Big)R\ln2$.}
\note{парадокс Гиббса: для \emph{одинаковых} газов формула дала бы $\Delta S>0$, хотя состояние не изменилось. Для тождественных газов $\Delta S=0$.}
\end{problem}

%-------------------------------------------------------------
\begin{problem}{2.11}{ex}{Какую максимальную работу можно получить, используя два тела с температурами $T_{10}$ и $T_{20}$ ($T_{10}>T_{20}$) как нагреватель и холодильник?
Теплоёмкости $C_1$ и $C_2$ постоянны. Найти конечную температуру $T$.}
\textbf{1.} Работа максимальна, если процесс обратимый (машина Карно с бесконечно малыми циклами); тогда энтропия системы не меняется:
\[C_1\ln\frac{T}{T_{10}}+C_2\ln\frac{T}{T_{20}}=0\ \Rightarrow\ T^{C_1+C_2}=T_{10}^{C_1}T_{20}^{C_2}.\]

\textbf{2. Работа} = убыль энергии тел: $A=C_1(T_{10}-T)+C_2(T_{20}-T)$.
\answer{$T=T_{10}^{C_1/(C_1+C_2)}\,T_{20}^{C_2/(C_1+C_2)}$,\; $A_{\max}=C_1T_{10}+C_2T_{20}-(C_1+C_2)T$.}
\end{problem}

%-------------------------------------------------------------
\begin{problem}{2.16}{dz}{Вычислить КПД цикла, изображённого на рис.~2.6 (диаграмма $T$--$S$). Известны $T_1$, $T_2$, $S_1$, $S_2$ и заштрихованные площади
$\Delta Q_1$ и $\Delta Q_2$.}
На рисунке цикл вписан в прямоугольник $T_2\le T\le T_1$, $S_1\le S\le S_2$; заштрихованные области — части прямоугольника, \emph{не} входящие в цикл:
$\Delta Q_1$ — у верхней стороны, $\Delta Q_2$ — у нижней.

\textbf{1.} На $T$--$S$ теплота равна площади под кривой. Верхняя ветвь (вправо, $\dd S>0$) — тепло получено. Она проходит ниже стороны $T_1$ на площадь $\Delta Q_1$:
$Q_1=T_1(S_2-S_1)-\Delta Q_1$.

\textbf{2.} Нижняя ветвь (влево) — тепло отдано. Она проходит выше стороны $T_2$ на площадь $\Delta Q_2$:
$|Q_2|=T_2(S_2-S_1)+\Delta Q_2$.

\textbf{3.} $A=Q_1-|Q_2|=(T_1-T_2)(S_2-S_1)-\Delta Q_1-\Delta Q_2$ (площадь цикла), $\eta=A/Q_1$.
\answer{$\eta=\dfrac{(T_1-T_2)(S_2-S_1)-\Delta Q_1-\Delta Q_2}{T_1(S_2-S_1)-\Delta Q_1}$.}
\end{problem}

%-------------------------------------------------------------
\begin{problem}{2.19}{dz}{Цикл состоит из изобары 1--2, изохоры 2--3 и адиабаты 3--1. Показать, что для любого вещества с постоянными молярными теплоёмкостями
$C_V$ и $C_p$ температуры связаны соотношением $T_2/T_3=(T_2/T_1)^\gamma$, $\gamma=C_p/C_V$.}
\textbf{1.} Энтропия — функция состояния, за цикл $\Delta S=0$. Изменения по участкам:
изобара $\Delta S_{12}=C_p\ln\frac{T_2}{T_1}$, изохора $\Delta S_{23}=C_V\ln\frac{T_3}{T_2}$, адиабата $\Delta S_{31}=0$.

\textbf{2.} $C_p\ln\frac{T_2}{T_1}+C_V\ln\frac{T_3}{T_2}=0$. Делим на $C_V$: $\gamma\ln\frac{T_2}{T_1}=\ln\frac{T_2}{T_3}$, т.е.\
$\big(\frac{T_2}{T_1}\big)^\gamma=\frac{T_2}{T_3}$. \hfill$\blacksquare$
\end{problem}

%-------------------------------------------------------------
\begin{problem}{2.20}{dz}{Обратимый цикл: изотермическое расширение, изобарическое сжатие, адиабатическое сжатие (рис.~2.8).
Определить КПД, если отношение максимальной и минимальной абсолютных температур равно $\alpha$. $C_V$, $C_p$ постоянны.}
\textbf{1. Температуры.} Изотерма 1--2 — при $T_1=T_{\max}$. При изобарном сжатии 2--3 газ охлаждается до $T_3=T_{\min}=T_1/\alpha$; адиабата 3--1 возвращает к $T_1$.

\textbf{2. Энтропия} (за цикл $\Delta S=0$, на адиабате $0$): $\Delta S_{12}=-\Delta S_{23}=-C_p\ln\frac{T_3}{T_1}=C_p\ln\alpha$.

\textbf{3. Теплоты.} Изотерма: $Q_{12}=T_1\Delta S_{12}=C_pT_1\ln\alpha>0$ (получено). Изобара: $Q_{23}=C_p(T_3-T_1)<0$ (отдано).
\[\eta=1-\frac{|Q_{23}|}{Q_{12}}=1-\frac{C_p(T_1-T_1/\alpha)}{C_pT_1\ln\alpha}=1-\frac{1-1/\alpha}{\ln\alpha}.\]
\answer{$\eta=1-\dfrac{\alpha-1}{\alpha\ln\alpha}$.}
\end{problem}

%-------------------------------------------------------------
\begin{problem}{2.26}{dz}{Теплоизолированный цилиндр разделён пополам пористой перегородкой. Слева один моль водорода при $T_1=270$\,К, справа — моль азота при $T_2=290$\,К.
Перегородка пропускает водород и не пропускает азот. Найти $\Delta S$ после установления равновесия. Теплоёмкостью цилиндра пренебречь.}
\textbf{1. Конечное состояние.} Водород заполняет весь цилиндр ($V\to2V$), азот остаётся в своей половине.
Оба газа двухатомные, $C_V=\frac52R$; из сохранения энергии $\frac52R(T-T_1)+\frac52R(T-T_2)=0\Rightarrow T=\frac{T_1+T_2}2=280$\,К.

\textbf{2. Энтропия каждого газа:}
H$_2$: $\frac52R\ln\frac T{T_1}+R\ln2$ (нагрев + удвоение объёма); \;N$_2$: $\frac52R\ln\frac T{T_2}$ (только охлаждение).

\textbf{3. Сумма и числа:} $\frac52R\ln\frac{T^2}{T_1T_2}=\frac52\cdot8{,}31\cdot\ln\frac{280^2}{270\cdot290}=0{,}03$; $R\ln2=5{,}76$.
\answer{$\Delta S=\frac52R\ln\dfrac{(T_1+T_2)^2}{4T_1T_2}+R\ln2\approx5{,}8$ Дж/К (в книге $5{,}77$).}
\end{problem}

%-------------------------------------------------------------
\begin{problem}{2.27}{dz}{Сосуд объёмом $2V$ с теплонепроницаемыми стенками разделён пополам теплонепроницаемым поршнем. В каждой части по молю идеального газа
($C_V$ постоянна), начальные температуры $T_1$ и $T_2$. Поршень отпускают, его колебания затухают из-за внутреннего трения в газе.
После остановки поршень делит объём в отношении $V_1/V_2=n$. Трение поршня о стенки не учитывать. Найти $\Delta S$ газа.}
\textbf{1. Конечные объёмы:} $V_1'+V_2'=2V$, $V_1'/V_2'=n$ $\Rightarrow V_1'=\frac{2nV}{1+n}$, $V_2'=\frac{2V}{1+n}$.

\textbf{2. Конечные температуры.} Вся система изолирована, работа поршня — внутренняя, поэтому полная энергия газа сохраняется:
$C_V(T_1'+T_2')=C_V(T_1+T_2)$. Поршень остановился — давления равны: $\frac{RT_1'}{V_1'}=\frac{RT_2'}{V_2'}\Rightarrow\frac{T_1'}{T_2'}=n$. Отсюда
\[T_1'=\frac{n(T_1+T_2)}{1+n},\qquad T_2'=\frac{T_1+T_2}{1+n}.\]

\textbf{3. Энтропия:} $\Delta S=\sum_i\Big[C_V\ln\frac{T_i'}{T_i}+R\ln\frac{V_i'}{V}\Big]=C_V\ln\frac{T_1'T_2'}{T_1T_2}+R\ln\frac{V_1'V_2'}{V^2}$.
\answer{$\Delta S=C_V\ln\dfrac{n(T_1+T_2)^2}{(1+n)^2T_1T_2}+R\ln\dfrac{4n}{(1+n)^2}$.}
\end{problem}

%-------------------------------------------------------------
\begin{problem}{2.28}{dz}{Моль одноатомного идеального газа находится под поршнем в адиабатически изолированном цилиндре. Масса груза на поршне,
определяющая давление газа, мгновенно увеличивается вдвое. Найти изменение энтропии газа после установления равновесия. Атмосферным давлением пренебречь.}
\textbf{1. Процесс необратимый:} газ сжимается под постоянным внешним давлением $2p_0$ (груз). Работа над газом $2p_0(V_0-V)$, $Q=0$:
\[\tfrac32R(T-T_0)=2p_0(V_0-V).\]

\textbf{2. Конечное состояние.} Подставляем $p_0V_0=RT_0$ и $2p_0V=RT$ (в конце давление газа $2p_0$):
$\frac32R(T-T_0)=2RT_0-RT\Rightarrow T=\frac75T_0$. Объём: $\frac V{V_0}=\frac{T/2p_0}{T_0/p_0}=\frac{T}{2T_0}=0{,}7$.

\textbf{3. Энтропия} (формула через $T$ и $p$): $\Delta S=C_p\ln\frac T{T_0}-R\ln\frac p{p_0}=\frac52R\ln1{,}4-R\ln2=8{,}31\cdot(0{,}841-0{,}693)$.
\answer{$\Delta S=\frac52R\ln1{,}4-R\ln2\approx1{,}23$ Дж/К.}
\note{в книге то же в виде $\frac32R\ln2+\frac52R\ln0{,}7$.}
\end{problem}

%-------------------------------------------------------------
\begin{problem}{2.31}{dz}{В двух теплоизолированных сосудах находится по одному молю разных идеальных одноатомных газов при одинаковом давлении $p$.
Температура газа в первом сосуде $T_1$, во втором — $T_2$. Насколько изменится энтропия системы, если сосуды соединить?}
\textbf{1. Конечное состояние.} Объёмы $V_i=\frac{RT_i}{p}$. Энергия сохраняется: $\frac32R(T-T_1)+\frac32R(T-T_2)=0\Rightarrow T=\frac{T_1+T_2}2$.
Каждый газ занимает весь объём $V_1+V_2=\frac{R(T_1+T_2)}{p}$.

\textbf{2. Энтропия каждого газа:} $\Delta S_i=\frac32R\ln\frac{T}{T_i}+R\ln\frac{V_1+V_2}{V_i}=\frac32R\ln\frac{T}{T_i}+R\ln\frac{T_1+T_2}{T_i}$.

\textbf{3. Складываем:} $\Delta S=\frac32R\ln\frac{(T_1+T_2)^2}{4T_1T_2}+R\ln\frac{(T_1+T_2)^2}{T_1T_2}$, а
$R\ln\frac{(T_1+T_2)^2}{T_1T_2}=R\ln\frac{(T_1+T_2)^2}{4T_1T_2}+2R\ln2$.
\answer{$\Delta S=\frac52R\ln\dfrac{(T_1+T_2)^2}{4T_1T_2}+2R\ln2$.}
\end{problem}

\section*{§3. Термодинамические потенциалы}
\addcontentsline{toc}{section}{§3. Термодинамические потенциалы}

\begin{basics}
\sub{Четыре потенциала и их «родные» переменные}
\begin{tabular}{@{}lll@{}}
$U(S,V)$ & внутренняя энергия & $\dd U=T\dd S-p\dd V$\\
$F(T,V)=U-TS$ & свободная энергия & $\dd F=-S\dd T-p\dd V$\\
$H(S,p)=U+pV$ & энтальпия & $\dd H=T\dd S+V\dd p$\\
$G(T,p)=U-TS+pV$ & потенциал Гиббса & $\dd G=-S\dd T+V\dd p$
\end{tabular}
\\Равенства — для обратимых процессов (для необратимых вместо $=$ стоит $\le$).
$U(S,V)$ — та же внутренняя энергия, что в §1, но записанная через $S$ и $V$: тогда $T=U'_S$, $p=-U'_V$.\\
Потенциалом функция является только в своих переменных: $F$ — в $(T,V)$, $G$ — в $(T,p)$ и т.д.

\sub{Первые производные (читаются прямо с дифференциалов)}
\[T=\Big(\frac{\partial U}{\partial S}\Big)_V=\Big(\frac{\partial H}{\partial S}\Big)_p,\quad
p=-\Big(\frac{\partial U}{\partial V}\Big)_S=-\Big(\frac{\partial F}{\partial V}\Big)_T,\]
\[S=-\Big(\frac{\partial F}{\partial T}\Big)_V=-\Big(\frac{\partial G}{\partial T}\Big)_p,\quad
V=\Big(\frac{\partial H}{\partial p}\Big)_S=\Big(\frac{\partial G}{\partial p}\Big)_T.\]
Зная один потенциал в своих переменных, получаем всё: например, из $F(T,V)$: $p=-F'_V$ — уравнение состояния,
$S=-F'_T$, $U=F+TS$.

\sub{Соотношения Максвелла (равенство смешанных производных)}
\[\Big(\frac{\partial S}{\partial V}\Big)_T=\Big(\frac{\partial p}{\partial T}\Big)_V\ (\text{из }F),\qquad
\Big(\frac{\partial S}{\partial p}\Big)_T=-\Big(\frac{\partial V}{\partial T}\Big)_p\ (\text{из }G),\]
\[\Big(\frac{\partial T}{\partial V}\Big)_S=-\Big(\frac{\partial p}{\partial S}\Big)_V\ (\text{из }U),\qquad
\Big(\frac{\partial T}{\partial p}\Big)_S=\Big(\frac{\partial V}{\partial S}\Big)_p\ (\text{из }H).\]
Самые нужные — первые две: они заменяют неизмеримую производную энтропии на производную из уравнения состояния.

\sub{Теплоёмкости и полезные формулы}
$\widetilde C$ — теплоёмкость всего тела (Дж/К), $C$ без тильды — молярная.
\[\widetilde C_V=T\Big(\frac{\partial S}{\partial T}\Big)_V,\qquad \widetilde C_p=T\Big(\frac{\partial S}{\partial T}\Big)_p,\]
\[\Big(\frac{\partial U}{\partial V}\Big)_T=T\Big(\frac{\partial p}{\partial T}\Big)_V-p,\qquad
\Big(\frac{\partial H}{\partial p}\Big)_T=V-T\Big(\frac{\partial V}{\partial T}\Big)_p,\]
\[\widetilde C_p-\widetilde C_V=T\Big(\frac{\partial p}{\partial T}\Big)_V\Big(\frac{\partial V}{\partial T}\Big)_p.\]
Уравнения Гиббса–Гельмгольца: $U=F-T\big(\frac{\partial F}{\partial T}\big)_V$, $H=G-T\big(\frac{\partial G}{\partial T}\big)_p$.

\sub{Математика: как переставлять производные}
\begin{itemize}
\item Тройное произведение: $\big(\frac{\partial x}{\partial y}\big)_z\big(\frac{\partial y}{\partial z}\big)_x\big(\frac{\partial z}{\partial x}\big)_y=-1$, т.е.
$\big(\frac{\partial x}{\partial y}\big)_z=-\dfrac{(\partial z/\partial y)_x}{(\partial z/\partial x)_y}$.
\item Обратная производная: $\big(\frac{\partial x}{\partial y}\big)_z=1\big/\big(\frac{\partial y}{\partial x}\big)_z$.
\item Смена постоянной переменной: $\big(\frac{\partial x}{\partial y}\big)_z=\big(\frac{\partial x}{\partial y}\big)_w+\big(\frac{\partial x}{\partial w}\big)_y\big(\frac{\partial w}{\partial y}\big)_z$.
\end{itemize}
\textbf{Как доказывать тождество:} (1) производную с $S=\mathrm{const}$ переписать через тройное произведение;
(2) $(\partial S/\partial T)$ заменить на $\widetilde C/T$; (3) $(\partial S/\partial V)_T$, $(\partial S/\partial p)_T$ — по Максвеллу.

\sub{Равновесие и работа}
Изолированная система: $S\to\max$. При $T,V=\mathrm{const}$: $F\to\min$; при $T,p=\mathrm{const}$: $G\to\min$.
Изотермический процесс: $A\le-\Delta F$ ($A_{\max}=-\Delta F$). Изобарический: $Q=\Delta H$.
\end{basics}

%-------------------------------------------------------------
\begin{problem}{3.1}{exdz}{Найти потенциал Гиббса $G$ и свободную энергию $F$ для $\nu$ молей идеального газа.}
\textbf{Найти:} $F(T,V)$ и $G(T,p)$.

\textbf{1. Исходные $U$ и $S$ идеального газа} ($C_V$ — молярная теплоёмкость, $S_0$ — константа):
\[U=\nu C_VT,\qquad S=\nu\,(C_V\ln T+R\ln V+S_0).\]

\textbf{2. Свободная энергия} $F=U-TS$:
\[F(T,V)=\nu\big[C_VT-T(C_V\ln T+R\ln V+S_0)\big]=\nu\big[C_VT(1-\ln T)-RT\ln V-TS_0\big].\]

\textbf{3. Потенциал Гиббса} $G=F+pV$, где $pV=\nu RT$. Его надо выразить через $T,p$: подставляем $V=\nu RT/p$, т.е.
$\ln V=\ln T-\ln p+\ln(\nu R)$:
\[G=\nu\big[C_VT-C_VT\ln T-RT\ln T+RT\ln p-RT\ln(\nu R)-TS_0+RT\big].\]
Собираем, учитывая $C_V+R=C_p$, и обозначаем $S_0'=S_0+R\ln(\nu R)$:
\[G(T,p)=\nu\big[C_pT(1-\ln T)+RT\ln p-TS_0'\big].\]
\textit{Проверка:} $V=(\partial G/\partial p)_T=\nu RT/p$ — верно.
\answer{$F=\nu[C_VT(1-\ln T)-RT\ln V-TS_0]$,\quad $G=\nu[C_pT(1-\ln T)+RT\ln p-TS_0']$.}
\end{problem}

%-------------------------------------------------------------
\begin{problem}{3.2}{ex}{Найти уравнение состояния одного моля газа, если его энтальпия $H=C_pp^{(\gamma-1)/\gamma}\exp[(S-S_0)/C_p]$, где
$C_p$ — молярная теплоёмкость при постоянном давлении, $S_0$ — константа, $\gamma=C_p/C_V$.}
\textbf{Найти:} связь $p,V,T$.

\textbf{1.} $H$ дана в своих переменных $(S,p)$, поэтому $T=(\partial H/\partial S)_p$, $V=(\partial H/\partial p)_S$:
\[T=p^{(\gamma-1)/\gamma}e^{(S-S_0)/C_p},\qquad
V=C_p\,\frac{\gamma-1}{\gamma}\,p^{-1/\gamma}e^{(S-S_0)/C_p}.\]

\textbf{2.} Делим, чтобы исключить $S$: $\dfrac TV=\dfrac{\gamma\,p}{C_p(\gamma-1)}=\dfrac{p}{C_p-C_V}=\dfrac pR$
(так как $C_p(\gamma-1)/\gamma=C_p-C_V=R$).
\answer{$pV=RT$ — идеальный газ.}
\end{problem}

%-------------------------------------------------------------
\begin{problem}{3.3}{exdz}{Показать, что теплоёмкость вещества при постоянном давлении $p$, объём $V$ которого линейно зависит от температуры
$T$, не зависит от величины $p$.}
\textbf{Найти (доказать):} $(\partial\widetilde C_p/\partial p)_T=0$, если $V=a(p)+b(p)T$.

\textbf{1.} По определению $\widetilde C_p=T(\partial S/\partial T)_p$. Дифференцируем по $p$ при $T=\mathrm{const}$ и меняем порядок дифференцирования:
\[\Big(\frac{\partial\widetilde C_p}{\partial p}\Big)_T=T\,\frac{\partial^2S}{\partial p\,\partial T}
=T\,\frac{\partial}{\partial T}\Big(\frac{\partial S}{\partial p}\Big)_T.\]

\textbf{2.} По Максвеллу $(\partial S/\partial p)_T=-(\partial V/\partial T)_p$:
\[\Big(\frac{\partial\widetilde C_p}{\partial p}\Big)_T=-T\Big(\frac{\partial^2V}{\partial T^2}\Big)_p.\]

\textbf{3.} $V$ линеен по $T$ $\Rightarrow$ $(\partial^2V/\partial T^2)_p=0$.
\answer{$(\partial\widetilde C_p/\partial p)_T=-T(\partial^2V/\partial T^2)_p=0$, т.е.\ $\widetilde C_p$ от $p$ не зависит.}
\end{problem}

%-------------------------------------------------------------
\begin{problem}{3.4}{ex}{Найти разность теплоёмкостей $\widetilde C_p$ и $\widetilde C_V$ для вещества, уравнение состояния которого
$V=V(p,T)$ известно.}
\textbf{Найти:} $\widetilde C_p-\widetilde C_V$ через производные $V(p,T)$.

\textbf{1. Из первого начала} (пример 1.7): $\widetilde C_p-\widetilde C_V=\big[(\partial U/\partial V)_T+p\big](\partial V/\partial T)_p$.

\textbf{2. Внутренняя энергия.} Из $\dd U=T\dd S-p\dd V$ и Максвелла $(\partial S/\partial V)_T=(\partial p/\partial T)_V$:
\[\Big(\frac{\partial U}{\partial V}\Big)_T+p=T\Big(\frac{\partial p}{\partial T}\Big)_V
\ \Rightarrow\ \widetilde C_p-\widetilde C_V=T\Big(\frac{\partial p}{\partial T}\Big)_V\Big(\frac{\partial V}{\partial T}\Big)_p.\]

\textbf{3. Через $V(p,T)$.} Тройное произведение: $(\partial p/\partial T)_V=-(\partial V/\partial T)_p\big/(\partial V/\partial p)_T$.
\answer{$\widetilde C_p-\widetilde C_V=-T\,\dfrac{(\partial V/\partial T)_p^2}{(\partial V/\partial p)_T}$\quad ($>0$, т.к.\ $(\partial V/\partial p)_T<0$).}
\note{для моля идеального газа получается $C_p-C_V=R$.}
\end{problem}

%-------------------------------------------------------------
\begin{problem}{3.5}{ex}{Объём однородного изотропного вещества, уравнение состояния которого $V=V(T,p)$ и теплоёмкость $C_p$ известны,
изменяется в ходе адиабатического квазистатического процесса. Вывести дифференциальное уравнение, связывающее давление и температуру
вещества. Решить его в случае одного моля одноатомного идеального газа. Начальные температура и давление равны $T_0$ и $p_0$.}
\textbf{Найти:} уравнение для $\dd T/\dd p$ при $S=\mathrm{const}$; $T(p)$ для одноатомного газа.

\textbf{1. Адиабата квазистатическая $\Rightarrow$ $S=\mathrm{const}$:}
$\dd S=\big(\frac{\partial S}{\partial T}\big)_p\dd T+\big(\frac{\partial S}{\partial p}\big)_T\dd p=0$.

\textbf{2. Подставляем} $(\partial S/\partial T)_p=\widetilde C_p/T$ и $(\partial S/\partial p)_T=-(\partial V/\partial T)_p$ (Максвелл):
\[\frac{\widetilde C_p}{T}\,\frac{\dd T}{\dd p}=\Big(\frac{\partial V}{\partial T}\Big)_p.\]

\textbf{3. Одноатомный газ:} $C_p=\frac52R$, $V=RT/p$, $(\partial V/\partial T)_p=R/p$:
\[\frac52\,\frac{\dd T}{T}=\frac{\dd p}{p}\ \Rightarrow\ \frac52\ln\frac T{T_0}=\ln\frac p{p_0}\ \Rightarrow\ T=T_0\Big(\frac p{p_0}\Big)^{2/5}.\]
\answer{$\dfrac{\widetilde C_p}{T}\dfrac{\dd T}{\dd p}=\Big(\dfrac{\partial V}{\partial T}\Big)_p$;\quad
$\Delta T=T_0\big[(p/p_0)^{2/5}-1\big]$.}
\end{problem}

%-------------------------------------------------------------
\begin{problem}{3.6}{ex}{Воду, находящуюся при температуре $T_0=273$\,К и давлении $p_0=100$\,атм, квазистатически адиабатически расширяют
до атмосферного давления $\tilde p$. Найти изменение температуры воды, если термический коэффициент её объёмного расширения
$\alpha=-6{,}4\cdot10^{-5}$\,К$^{-1}$, плотность $\rho=10^3$\,кг/м$^3$ остаётся постоянной, а удельная теплоёмкость $c_p=4{,}2$\,кДж/(кг$\cdot$К).}
\textbf{Найти:} $T-T_0$.

\textbf{1. Уравнение адиабаты из 3.5} на 1 кг: $\frac{c_p}{T}\dd T=\big(\frac{\partial v}{\partial T}\big)_p\dd p$, где удельный объём
$v=1/\rho$ и $(\partial v/\partial T)_p=\alpha v=\alpha/\rho$:
\[\frac{\dd T}{T}=\frac{\alpha}{\rho c_p}\dd p\ \Rightarrow\ T=T_0\exp\frac{\alpha(\tilde p-p_0)}{\rho c_p}\approx T_0\Big(1+\frac{\alpha(\tilde p-p_0)}{\rho c_p}\Big).\]

\textbf{2. Числа:} $\tilde p-p_0=-99$\,атм $=-1{,}0\cdot10^7$\,Па;
\[\frac{\alpha(\tilde p-p_0)}{\rho c_p}=\frac{(-6{,}4\cdot10^{-5})(-1{,}0\cdot10^7)}{10^3\cdot4{,}2\cdot10^3}\approx1{,}5\cdot10^{-4},\quad
T-T_0\approx273\cdot1{,}5\cdot10^{-4}.\]
\answer{$T-T_0\approx+0{,}04$\,К — вода при расширении нагревается (т.к.\ $\alpha<0$ при 0–4\,$^\circ$C).}
\end{problem}

%-------------------------------------------------------------
\begin{problem}{3.7}{ex}{Процесс Джоуля–Томсона: газ из части 1 цилиндра медленно продавливается через пористую перегородку в часть 2;
поршни поддерживают постоянные давления $p_1$ и $p_2$ ($p_1>p_2$), цилиндр теплоизолирован. Начальная температура $T_0$.
Определить конечную температуру $T_1$.}
\textbf{Найти:} $T_1$ через уравнение состояния и $\widetilde C_p$.

\textbf{1. Что сохраняется.} Слева над газом совершают работу $p_1V_1$, справа газ совершает $p_2V_2$, теплообмена нет:
$U_2-U_1=p_1V_1-p_2V_2$ $\Rightarrow$ $U_1+p_1V_1=U_2+p_2V_2$, т.е.\ $H=\mathrm{const}$.

\textbf{2. $\dd H=0$ в переменных $T,p$:} $\dd H=\big(\frac{\partial H}{\partial T}\big)_p\dd T+\big(\frac{\partial H}{\partial p}\big)_T\dd p=0$, где
\[\Big(\frac{\partial H}{\partial T}\Big)_p=T\Big(\frac{\partial S}{\partial T}\Big)_p=\widetilde C_p,\qquad
\Big(\frac{\partial H}{\partial p}\Big)_T=T\Big(\frac{\partial S}{\partial p}\Big)_T+V=V-T\Big(\frac{\partial V}{\partial T}\Big)_p.\]

\textbf{3. Отсюда коэффициент Джоуля–Томсона и интеграл:}
\[\Big(\frac{\partial T}{\partial p}\Big)_H=\frac{T(\partial V/\partial T)_p-V}{\widetilde C_p},\qquad
T_1=T_0+\int_{p_1}^{p_2}\frac{T(\partial V/\partial T)_p-V}{\widetilde C_p}\,\dd p.\]
\answer{$T_1=T_0+\displaystyle\int_{p_1}^{p_2}\frac{1}{\widetilde C_p}\Big[T\Big(\frac{\partial V}{\partial T}\Big)_p-V\Big]\dd p$.}
\note{для идеального газа $T(\partial V/\partial T)_p=V$ $\Rightarrow$ $T_1=T_0$ (задача 3.22).}
\end{problem}

%-------------------------------------------------------------
\begin{problem}{3.8}{exdz}{Доказать соотношение $(\partial T/\partial V)_S=-T(\partial p/\partial T)_V/\widetilde C_V$.}
\textbf{1. Тройное произведение} для $T,V,S$ (уходим от $S=\mathrm{const}$):
\[\Big(\frac{\partial T}{\partial V}\Big)_S=-\frac{(\partial S/\partial V)_T}{(\partial S/\partial T)_V}.\]

\textbf{2. Знаменатель:} $(\partial S/\partial T)_V=\widetilde C_V/T$.
\textbf{Числитель} по Максвеллу: $(\partial S/\partial V)_T=(\partial p/\partial T)_V$.

\textbf{3. Подставляем:} $\Big(\dfrac{\partial T}{\partial V}\Big)_S=-\dfrac{(\partial p/\partial T)_V}{\widetilde C_V/T}=-\dfrac{T(\partial p/\partial T)_V}{\widetilde C_V}$.
\answer{доказано.}
\end{problem}

%-------------------------------------------------------------
\begin{problem}{3.9}{ex}{Термодинамическая система находится в термостате при постоянном давлении. Найти условия существования устойчивого
равновесия такой системы.}
\textbf{Найти:} неравенства на производные, обеспечивающие устойчивость.

\textbf{1. Какой потенциал минимален.} $T,p$ заданы $\Rightarrow$ в равновесии $G=U(S,V)-TS+pV\to\min$ по $S$ и $V$.

\textbf{2. Первая вариация $=0$:}
$\delta G=\big[(\partial U/\partial V)_S+p\big]\delta V+\big[(\partial U/\partial S)_V-T\big]\delta S=0$
$\Rightarrow$ температура и давление системы равны внешним.

\textbf{3. Вторая вариация $>0$} (положительная квадратичная форма вторых производных $U$):
\[\delta^2G=\tfrac12U_{VV}\delta V^2+U_{VS}\delta V\delta S+\tfrac12U_{SS}\delta S^2>0\]
\[\iff\quad U_{SS}>0,\qquad U_{VV}U_{SS}-U_{VS}^2>0.\]

\textbf{4. Смысл.} $U_{SS}=(\partial T/\partial S)_V=T/\widetilde C_V>0$ $\Rightarrow$ $\widetilde C_V>0$.
Второе условие (переход к переменным $T,V$) сводится к $-T(\partial p/\partial V)_T/\widetilde C_V>0$.
\answer{$\widetilde C_V>0$ (тепловая устойчивость) и $(\partial p/\partial V)_T<0$ (механическая).}
\end{problem}

%-------------------------------------------------------------
\begin{problem}{3.10}{ex}{В объёмах $V_1=3$\,м$^3$ и $V_2=2$\,м$^3$ при $T_0=300$\,К находятся по $\nu=0{,}5$ моля кислорода и азота
соответственно. Оценить максимальную работу, которая может быть произведена за счёт изотермического смешивания этих газов в объёме
$V_1+V_2$. Газы считать идеальными.}
\textbf{Найти:} $A_{\max}$.

\textbf{1. Изотермический процесс $\Rightarrow$ $A_{\max}=-\Delta F$.} Газы идеальные, каждый просто расширяется на весь объём
$V=V_1+V_2$. Из 3.1 при $T=\mathrm{const}$ меняется только слагаемое $-\nu RT\ln V$:
\[\Delta F=-\nu RT_0\ln\frac V{V_1}-\nu RT_0\ln\frac V{V_2}=-\nu RT_0\ln\frac{(V_1+V_2)^2}{V_1V_2}.\]

\textbf{2. Числа:} $\frac{(V_1+V_2)^2}{V_1V_2}=\frac{25}6$, $\ln\frac{25}6=1{,}43$; $\nu RT_0=0{,}5\cdot8{,}31\cdot300=1247$\,Дж;
$A_{\max}=1247\cdot1{,}43$.
\answer{$A_{\max}=\nu RT_0\ln\dfrac{(V_1+V_2)^2}{V_1V_2}\approx1{,}8$\,кДж.}
\end{problem}

%-------------------------------------------------------------
\begin{problem}{3.16}{dz}{Доказать соотношения $U=F-T(\partial F/\partial T)_V$ и $H=G-T(\partial G/\partial T)_p$.}
\textbf{1.} По определению $F=U-TS$ $\Rightarrow$ $U=F+TS$. Из $\dd F=-S\dd T-p\dd V$: $S=-(\partial F/\partial T)_V$.
Подставляем: $U=F-T(\partial F/\partial T)_V$.

\textbf{2.} По определению $G=H-TS$ $\Rightarrow$ $H=G+TS$. Из $\dd G=-S\dd T+V\dd p$: $S=-(\partial G/\partial T)_p$.
Подставляем: $H=G-T(\partial G/\partial T)_p$.
\answer{доказано (уравнения Гиббса–Гельмгольца).}
\end{problem}

%-------------------------------------------------------------
\begin{problem}{3.17}{dz}{Доказать соотношение $(\partial T/\partial p)_S=T(\partial V/\partial T)_p/\widetilde C_p$.}
\textbf{1. Тройное произведение} для $T,p,S$:
\[\Big(\frac{\partial T}{\partial p}\Big)_S=-\frac{(\partial S/\partial p)_T}{(\partial S/\partial T)_p}.\]

\textbf{2.} $(\partial S/\partial T)_p=\widetilde C_p/T$; по Максвеллу (из $\dd G$) $(\partial S/\partial p)_T=-(\partial V/\partial T)_p$.

\textbf{3. Подставляем:} $\Big(\dfrac{\partial T}{\partial p}\Big)_S=\dfrac{(\partial V/\partial T)_p}{\widetilde C_p/T}=\dfrac{T(\partial V/\partial T)_p}{\widetilde C_p}$.
\answer{доказано (аналог 3.8 с заменой $V\to p$, $\widetilde C_V\to\widetilde C_p$).}
\end{problem}

%-------------------------------------------------------------
\begin{problem}{3.18}{dz}{Доказать соотношения $(\partial U/\partial V)_T=-\widetilde C_V(\partial T/\partial V)_S-p$ и
$(\partial H/\partial p)_T=-\widetilde C_p(\partial T/\partial p)_S+V$.}
\textbf{1. Первое.} Из $\dd U=T\dd S-p\dd V$ при $T=\mathrm{const}$ и Максвелла $(\partial S/\partial V)_T=(\partial p/\partial T)_V$:
\[\Big(\frac{\partial U}{\partial V}\Big)_T=T\Big(\frac{\partial S}{\partial V}\Big)_T-p=T\Big(\frac{\partial p}{\partial T}\Big)_V-p.\]
По 3.8 $T(\partial p/\partial T)_V=-\widetilde C_V(\partial T/\partial V)_S$ $\Rightarrow$
$(\partial U/\partial V)_T=-\widetilde C_V(\partial T/\partial V)_S-p$.

\textbf{2. Второе.} Из $\dd H=T\dd S+V\dd p$ при $T=\mathrm{const}$ и Максвелла $(\partial S/\partial p)_T=-(\partial V/\partial T)_p$:
\[\Big(\frac{\partial H}{\partial p}\Big)_T=T\Big(\frac{\partial S}{\partial p}\Big)_T+V=-T\Big(\frac{\partial V}{\partial T}\Big)_p+V.\]
По 3.17 $T(\partial V/\partial T)_p=\widetilde C_p(\partial T/\partial p)_S$ $\Rightarrow$
$(\partial H/\partial p)_T=-\widetilde C_p(\partial T/\partial p)_S+V$.
\answer{доказано.}
\end{problem}

%-------------------------------------------------------------
\begin{problem}{3.20}{dz}{Свободная энергия $F$ одного моля некоторого вещества задаётся выражением $F=-0{,}5RT\ln(AT^3V^2)$,
где $A$ — некоторая константа. Найти молярную теплоёмкость $C_p$ этого вещества.}
\textbf{Найти:} $C_p=T(\partial S/\partial T)_p$.

\textbf{1. Энтропия и уравнение состояния из $F(T,V)$.} Запишем $F=-\frac12RT(\ln A+3\ln T+2\ln V)$:
\[S=-\Big(\frac{\partial F}{\partial T}\Big)_V=\tfrac12R\ln(AT^3V^2)+\tfrac32R,\qquad
p=-\Big(\frac{\partial F}{\partial V}\Big)_T=\frac{RT}{V}.\]

\textbf{2. Энтропия в переменных $T,p$.} Подставляем $V=RT/p$: $AT^3V^2=AR^2T^5/p^2$,
\[S=\tfrac52R\ln T-R\ln p+\mathrm{const}.\]

\textbf{3. Теплоёмкость:} $C_p=T\big(\frac{\partial S}{\partial T}\big)_p=T\cdot\frac{5R}{2T}$.
\answer{$C_p=\frac52R$ (вещество — одноатомный идеальный газ: $pV=RT$, $U=F+TS=\frac32RT$).}
\end{problem}

\section*{§8. Распределение Максвелла}
\addcontentsline{toc}{section}{§8. Распределение Максвелла}

\begin{basics}
\sub{Что такое плотность вероятности $w$}
$w(x)\,\dd x$ — вероятность того, что величина у одной молекулы попадёт в интервал $(x,x+\dd x)$
(= средняя доля молекул с таким $x$). Число молекул: $\dd N=Nw(x)\dd x$.
Нормировка $\int w\,\dd x=1$. Среднее: $\avg{f}=\int f(x)w(x)\dd x$; дисперсия $\sigma^2=\avg{x^2}-\avg x^2$;
относительная флуктуация $\delta=\sigma/\avg x$.
Независимые величины: $w(x,y)=w(x)w(y)$.

\sub{Проекция скорости — гаусс}
\[w(v_x)=\sqrt{\frac{m}{2\pi kT}}\exp\Big(-\frac{mv_x^2}{2kT}\Big),\qquad \avg{v_x}=0,\quad\avg{v_x^2}=\frac{kT}{m}.\]
$m$ — масса \emph{одной} молекулы ($m=\mu/N_A$, $\frac{kT}{m}=\frac{RT}{\mu}$), $T$ — температура газа.
Проекции $v_x,v_y,v_z$ независимы.

\sub{Вектор скорости}
\[w(\vec v)=w(v_x)w(v_y)w(v_z)=\Big(\frac{m}{2\pi kT}\Big)^{3/2}\exp\Big(-\frac{mv^2}{2kT}\Big)\]
— плотность на единицу объёма \emph{пространства скоростей} $\dd^3v=\dd v_x\dd v_y\dd v_z$. Зависит только от $|\vec v|$: все направления равновероятны.

\sub{Модуль скорости (распределение Максвелла)}
Слой $(v,v+\dd v)$ в пространстве скоростей имеет объём $4\pi v^2\dd v$:
\[w(v)=4\pi\Big(\frac{m}{2\pi kT}\Big)^{3/2}v^2\exp\Big(-\frac{mv^2}{2kT}\Big).\]
Характерные скорости:
\begin{align*}
\vp&=\sqrt{2kT/m} &&\text{наивероятнейшая (максимум $w$)},\\
\vav&=\sqrt{8kT/\pi m} &&\text{средняя},\\
v_{\text{кв}}&=\sqrt{\avg{v^2}}=\sqrt{3kT/m} &&\text{среднеквадратичная};\qquad \vp<\vav<v_{\text{кв}}.
\end{align*}
Средняя кинетическая энергия $\avg E=\frac m2\avg{v^2}=\frac32kT$.

\sub{Замена переменной}
Если $y=y(x)$ монотонна: $w(y)=w(x(y))\,\Big|\dfrac{\dd x}{\dd y}\Big|$ (вероятность в соответствующих интервалах одна и та же).
Для нескольких переменных — на якобиан. Максимум $w(y)$ \emph{не} совпадает с $y(x^*)$.

\sub{Поток молекул на стенку (эффузия)}
Число ударов о единицу площади за единицу времени:
$j=n\int_0^\infty v_xw(v_x)\dd v_x=\frac14n\vav$.
Через малое отверстие площади $S$ вылетает $\frac14nS\vav$ молекул в секунду.

\sub{Гауссовы интегралы}
\[I_k=\int_0^\infty x^{k}e^{-ax^2}\dd x=\frac{\Gamma\big(\frac{k+1}2\big)}{2a^{(k+1)/2}},\qquad I_{k+2}=\frac{k+1}{2a}I_k:\]
\[I_0=\frac12\sqrt{\frac\pi a},\ I_1=\frac1{2a},\ I_2=\frac{\sqrt\pi}{4a^{3/2}},\ I_3=\frac1{2a^2},\ I_4=\frac{3\sqrt\pi}{8a^{5/2}},\ I_5=\frac1{a^3}.\]
$\int_{-\infty}^\infty=2I_k$ для чётных $k$ и $0$ для нечётных.
Функция ошибок $\erf x=\frac2{\sqrt\pi}\int_0^xe^{-t^2}\dd t$, $\erf\infty=1$.
\end{basics}

%-------------------------------------------------------------
\begin{problem}{8.1}{ex}{Оценить отношение числа молекул азота, проекция скорости которых на произвольную ось лежит в интервале
$(v_1-\Delta v_1,\,v_1+\Delta v_1)$, к числу молекул с проекцией в интервале $(v_2-\Delta v_2,\,v_2+\Delta v_2)$, при $T=300$\,К.
$\mu=28$\,г/моль, $v_1=300$\,м/с, $\Delta v_1=0{,}31$\,м/с, $v_2=500$\,м/с, $\Delta v_2=0{,}51$\,м/с.}
\textbf{1.} Интервалы узкие ($\Delta v\ll v$), поэтому число молекул в интервале $\approx$ плотность $\times$ ширина:
$\Delta N_i=N\,w(v_i)\cdot2\Delta v_i$, где $w(v_x)=\sqrt{\frac m{2\pi kT}}e^{-mv_x^2/2kT}$.

\textbf{2.} В отношении $N$ и предэкспоненциальный множитель сокращаются:
\[\frac{\Delta N_1}{\Delta N_2}=\frac{\Delta v_1}{\Delta v_2}\exp\Big[\frac{m(v_2^2-v_1^2)}{2kT}\Big].\]

\textbf{3. Числа:} $m=\mu/N_A=4{,}65\cdot10^{-26}$\,кг;\;
$\frac{m(v_2^2-v_1^2)}{2kT}=\frac{4{,}65\cdot10^{-26}\cdot1{,}6\cdot10^5}{2\cdot1{,}38\cdot10^{-23}\cdot300}=0{,}90$;\;
$\frac{0{,}31}{0{,}51}e^{0{,}90}=0{,}61\cdot2{,}46$.
\answer{$\Delta N_1/\Delta N_2\approx1{,}5$.}
\end{problem}

%-------------------------------------------------------------
\begin{problem}{8.2}{exdz}{Используя распределение Максвелла, найти плотность вероятности распределения молекул по кинетическим энергиям $w(E)$.}
\textbf{1. Замена переменной.} $E=\frac{mv^2}{2}\Rightarrow v=\sqrt{\frac{2E}{m}}$, $\frac{\dd v}{\dd E}=\frac{1}{\sqrt{2mE}}$.
Вероятность попасть в соответствующие интервалы одна и та же: $w(E)\dd E=w(v)\dd v$, т.е.\ $w(E)=w(v)\frac{\dd v}{\dd E}$.

\textbf{2. Подставляем} $w(v)=4\pi\big(\frac m{2\pi kT}\big)^{3/2}v^2e^{-mv^2/2kT}$, $v^2=\frac{2E}m$:
\[w(E)=4\pi\Big(\frac m{2\pi kT}\Big)^{3/2}\frac{2E}{m}e^{-E/kT}\frac{1}{\sqrt{2mE}}=\frac{2}{\sqrt\pi}(kT)^{-3/2}\sqrt E\,e^{-E/kT}.\]

\textbf{3. Характерные значения.} Максимум: $\frac{\dd}{\dd E}\big(\sqrt Ee^{-E/kT}\big)=\big(\frac{1}{2\sqrt E}-\frac{\sqrt E}{kT}\big)e^{-E/kT}=0\Rightarrow E^*=\frac{kT}2$.
Среднее: $\avg E=\frac m2\avg{v^2}=\frac m2\cdot\frac{3kT}m=\frac32kT$.
\answer{$w(E)=\dfrac{2}{\sqrt\pi}(kT)^{-3/2}\sqrt E\,e^{-E/kT}$,\; $E^*=\dfrac{kT}2$,\; $\avg E=\dfrac32kT$.}
\note{$E^*\ne\frac{m\vp^2}{2}=kT$ — наиболее вероятные значения при замене переменной не пересчитываются.}
\end{problem}

%-------------------------------------------------------------
\begin{problem}{8.3}{exdz}{Для идеального газа из $N$ частиц найти относительные флуктуации энергии одной молекулы $\delta(E)$ и всего газа $\delta(E_N)$.}
\textbf{Найти:} $\delta=\sigma/\text{среднее}$, где $\sigma^2=\avg{E^2}-\avg E^2$.

\textbf{1. Средние по $w(E)$ из 8.2.} Замена $E=kTt$ сводит интегралы к $\int_0^\infty t^{k}e^{-t}\dd t=\Gamma(k+1)$:
\[\avg E=\frac{2kT}{\sqrt\pi}\Gamma\big(\tfrac52\big)=\frac{2kT}{\sqrt\pi}\cdot\frac{3\sqrt\pi}{4}=\frac32kT,\]
\[\avg{E^2}=\frac{2(kT)^2}{\sqrt\pi}\Gamma\big(\tfrac72\big)=\frac{2(kT)^2}{\sqrt\pi}\cdot\frac{15\sqrt\pi}{8}=\frac{15}4(kT)^2.\]

\textbf{2. Одна молекула:}
\[\sigma^2=\frac{15}4(kT)^2-\frac94(kT)^2=\frac32(kT)^2,\qquad \delta(E)=\frac{\sqrt{3/2}\,kT}{\frac32kT}=\sqrt{\frac23}.\]

\textbf{3. Весь газ.} $E_N=E_1+\dots+E_N$, молекулы независимы: средние и дисперсии складываются,
$\avg{E_N}=N\avg E$, $\sigma_N^2=N\sigma^2$, поэтому $\delta(E_N)=\frac{\sqrt N\sigma}{N\avg E}=\frac{\delta(E)}{\sqrt N}$.
\answer{$\delta(E)=\sqrt{2/3}\approx0{,}82$,\; $\delta(E_N)=\sqrt{2/(3N)}$.}
\end{problem}

%-------------------------------------------------------------
\begin{problem}{8.4}{ex}{Найти вероятность того, что кинетическая энергия молекулы идеального газа не превышает $E_0$. Температура $T$.}
\textbf{1.} $P(E_0)=\int_0^{E_0}w(E)\dd E$ с $w(E)$ из 8.2. Замена $E=kTt^2$ ($\dd E=2kTt\,\dd t$), $x_0=\sqrt{E_0/kT}$:
\[P=\frac{2}{\sqrt\pi}\int_0^{x_0}t\,e^{-t^2}\cdot2t\,\dd t=\frac4{\sqrt\pi}\int_0^{x_0}t^2e^{-t^2}\dd t.\]

\textbf{2. По частям} ($u=t$, $\dd v=te^{-t^2}\dd t$, $v=-\frac12e^{-t^2}$):
$\int_0^{x_0}t^2e^{-t^2}\dd t=-\frac{x_0}2e^{-x_0^2}+\frac12\int_0^{x_0}e^{-t^2}\dd t=-\frac{x_0}2e^{-x_0^2}+\frac{\sqrt\pi}4\erf x_0$.
\answer{$P(E_0)=\erf\sqrt{y}-2\sqrt{y/\pi}\;e^{-y}$, где $y=E_0/kT$.}
\note{энергию меньше $E^*=kT/2$ имеют около 20\% молекул, меньше $\avg E=\frac32kT$ — около 60\%.}
\end{problem}

%-------------------------------------------------------------
\begin{problem}{8.5}{ex}{Жидкость при температуре $T$ разлита тонким слоем; движение молекул плёнки двумерно, проекции скоростей независимы
и распределены по Гауссу. Найти $w(E)$, наиболее вероятные и средние значения энергии и модуля импульса. Масса молекулы $m$.}
\textbf{1. Распределение по модулю импульса.} $w(p_x,p_y)=\frac{1}{2\pi mkT}e^{-(p_x^2+p_y^2)/2mkT}$. В полярных координатах
элемент площади $p\,\dd p\,\dd\varphi$, интеграл по $\varphi$ даёт $2\pi$:
\[w(p)=\frac{p}{mkT}e^{-p^2/2mkT}.\]
Максимум: $p^*=\sqrt{mkT}$. Среднее: $\bar p=\frac1{mkT}\int_0^\infty p^2e^{-p^2/2mkT}\dd p=\sqrt{\frac{\pi mkT}2}$.

\textbf{2. По энергии.} $E=\frac{p^2}{2m}$, $\frac{\dd p}{\dd E}=\frac mp$:\; $w(E)=w(p)\frac mp=\frac1{kT}e^{-E/kT}$.
Это убывающая экспонента: $E^*=0$, $\avg E=kT$.
\answer{$w(E)=\dfrac1{kT}e^{-E/kT}$;\; $E^*=0$, $\avg E=kT$;\; $p^*=\sqrt{mkT}$, $\bar p=\sqrt{\pi mkT/2}$.}
\end{problem}

%-------------------------------------------------------------
\begin{problem}{8.6}{exdz}{Найти при температуре $T$ среднее значение модуля скорости относительного движения двух молекул идеального газа. Масса молекулы $m$.}
\textbf{1. Совместное распределение.} Молекулы независимы: $w(\vec v_1,\vec v_2)\propto e^{-m(v_1^2+v_2^2)/2kT}$.

\textbf{2. Новые переменные:} скорость центра масс $\vec V=\frac{\vec v_1+\vec v_2}2$ и относительная $\vec u=\vec v_1-\vec v_2$.
Тогда $v_1^2+v_2^2=2V^2+\frac{u^2}{2}$, и
\[w\propto e^{-2mV^2/2kT}\cdot e^{-(m/2)u^2/2kT}.\]
Распределение по $\vec u$ — максвелловское, но с массой $m/2$ вместо $m$.

\textbf{3. Средний модуль} — формула $\vav=\sqrt{8kT/\pi m}$ с заменой $m\to m/2$.
\answer{$\avg{|\vec v_1-\vec v_2|}=\sqrt{\dfrac{16kT}{\pi m}}=4\sqrt{\dfrac{kT}{\pi m}}=\sqrt2\,\vav$.}
\end{problem}

%-------------------------------------------------------------
\begin{problem}{8.7}{exdz}{Найти среднее число молекул, ударяющихся за время $\tau$ о плоскую стенку площадью $S$. Температура $T$, масса молекулы $m$, концентрация $n$.}
\textbf{1.} Ось $x$ — перпендикулярно стенке, к ней. За время $\tau$ о стенку ударятся молекулы с проекцией $v_x>0$,
находящиеся в цилиндре с основанием $S$ и высотой $v_x\tau$. Таких молекул с проекцией в $(v_x,v_x+\dd v_x)$:
$\dd N=n\,S\,v_x\tau\cdot w(v_x)\dd v_x$.

\textbf{2. Суммируем по всем $v_x>0$:}
\[N=nS\tau\sqrt{\frac{m}{2\pi kT}}\int_0^\infty v_xe^{-mv_x^2/2kT}\dd v_x=nS\tau\sqrt{\frac{m}{2\pi kT}}\cdot\frac{kT}{m}=nS\tau\sqrt{\frac{kT}{2\pi m}}.\]
\answer{$N=nS\tau\sqrt{\dfrac{kT}{2\pi m}}=\dfrac14nS\tau\,\vav$.}
\end{problem}

%-------------------------------------------------------------
\begin{problem}{8.8}{exdz}{Считая концентрацию $n$ практически постоянной, найти средние модуль скорости и энергию молекул, вылетающих из малого
отверстия площадью $S$ в плоской стенке сосуда. Температура газа $T$, масса молекулы $m$.}
\textbf{1. Кто вылетает.} За время $\tau$ из молекул со скоростью $v$ под углом $\theta$ к нормали вылетят те, что в наклонном цилиндре объёмом
$Sv\tau\cos\theta$. Значит, число вылетевших со скоростью $v$ пропорционально $v\cdot w(v)$: быстрые молекулы вылетают чаще.

\textbf{2. Средние по вылетевшим} — усредняем с весом $v\,w(v)$:
\[\avg v_{\text{выл}}=\frac{\int v\cdot v\,w\,\dd v}{\int v\,w\,\dd v}=\frac{\avg{v^2}}{\vav},\qquad
\avg E_{\text{выл}}=\frac m2\frac{\int v^2\cdot v\,w\,\dd v}{\int v\,w\,\dd v}=\frac m2\frac{\avg{v^3}}{\vav}.\]

\textbf{3. Подставляем} $\avg{v^2}=\frac{3kT}m$, $\vav=\sqrt{\frac{8kT}{\pi m}}$, $\avg{v^3}=\frac{4}{\sqrt\pi}\big(\frac{2kT}m\big)^{3/2}$
(гауссов интеграл $I_5$):
\[\avg v_{\text{выл}}=\frac{3kT/m}{\sqrt{8kT/\pi m}}=3\sqrt{\frac{\pi kT}{8m}},\qquad
\frac{\avg{v^3}}{\vav}=\frac{4kT}{m}\ \Rightarrow\ \avg E_{\text{выл}}=2kT.\]
\answer{$\avg v=3\sqrt{\dfrac{\pi kT}{8m}}=\dfrac{3\pi}8\vav$,\; $\avg E=2kT$ (больше, чем $\frac32kT$ в сосуде).}
\end{problem}

%-------------------------------------------------------------
\begin{problem}{8.9}{exdz}{Тонкостенный пустой сосуд объёмом $V$ погружён в идеальный газ, занимающий гораздо больший объём; температура $T$ и
концентрация $n_0$ газа постоянны. В стенке сосуда образовалось малое отверстие площадью $S$. Найти зависимость концентрации газа в сосуде
от времени. Средний модуль скорости молекул внутри и снаружи одинаков и равен $\vav$.}
\textbf{1. Потоки через отверстие} (8.7): внутрь за время $\dd t$ влетает $\frac14n_0\vav S\,\dd t$, наружу вылетает $\frac14n\vav S\,\dd t$.

\textbf{2. Баланс числа молекул в сосуде} ($N=nV$):
\[V\frac{\dd n}{\dd t}=\frac14S\vav(n_0-n),\qquad n(0)=0.\]

\textbf{3. Решаем} разделением переменных: $\frac{\dd n}{n_0-n}=\frac{\dd t}{\tau}$, $\tau=\frac{4V}{S\vav}$ $\Rightarrow$
$-\ln\frac{n_0-n}{n_0}=\frac t\tau$.
\answer{$n(t)=n_0\big(1-e^{-t/\tau}\big)$,\; $\tau=\dfrac{4V}{S\vav}$.}
\note{пример: He, 300\,К ($\vav\approx1260$\,м/с), $V=1$\,л, $S=0{,}01$\,мм$^2$ $\Rightarrow\tau\approx320$\,с.}
\end{problem}

%-------------------------------------------------------------
\begin{problem}{8.10}{exdz}{Невесомый поршень площадью $S$ удерживается в пустом цилиндре пружиной жёсткостью $\kappa$. В начале поршень у левой стенки,
пружина не деформирована. В левой стенке делают малое отверстие площадью $s_0$. Через какое время $\tau$ давление на поршень станет
равным половине атмосферного $p_0$? $T=\text{const}$, воздух — идеальный газ, сжатие пружины квазистатическое, $\vav$ одинакова внутри и снаружи.}
\textbf{1. Связь числа молекул с давлением.} Поршень сдвинулся на $x$, газ занимает объём $Sx$. Равновесие поршня: $\kappa x=pS\Rightarrow x=\frac{pS}{\kappa}$.
Уравнение состояния: $NkT=p\cdot Sx=\frac{p^2S^2}{\kappa}$, откуда $kT\,\dd N=\frac{2S^2}{\kappa}p\,\dd p$.

\textbf{2. Поток через отверстие} (как в 8.9), концентрации $n=\frac p{kT}$:
$\frac{\dd N}{\dd t}=\frac14s_0\vav\frac{p_0-p}{kT}$.

\textbf{3. Приравниваем} $\dd N$ из п.\,1 и п.\,2 и разделяем переменные:
\[\dd t=\frac{8S^2}{\kappa s_0\vav}\cdot\frac{p\,\dd p}{p_0-p}.\]

\textbf{4. Интегрируем} от $p=0$ до $p_0/2$; $\frac{p}{p_0-p}=-1+\frac{p_0}{p_0-p}$:
\[\int_0^{p_0/2}\frac{p\,\dd p}{p_0-p}=\Big[-p-p_0\ln(p_0-p)\Big]_0^{p_0/2}=p_0\Big(\ln2-\frac12\Big).\]
\answer{$\tau=\dfrac{8p_0S^2}{\kappa s_0\vav}\Big(\ln2-\dfrac12\Big)$.}
\end{problem}

%-------------------------------------------------------------
\begin{problem}{8.11}{ex}{Найти давление равновесного излучения на идеально отражающие стенки, считая его газом фотонов, движущихся в произвольных
направлениях с одинаковым по модулю импульсом $\hbar\omega/c$. Концентрация фотонов $n$.}
\textbf{1. Сколько фотонов ударяется.} Все летят со скоростью $c$, направления равновероятны: доля с углом к нормали в $(\theta,\theta+\dd\theta)$
равна $\frac12\sin\theta\,\dd\theta$. Как в 8.7, о площадку $S$ за время $\tau$ ударится $\dd N=nSc\tau\cos\theta\cdot\frac12\sin\theta\,\dd\theta$.

\textbf{2. Импульс.} При зеркальном отражении фотон передаёт стенке $2\frac{\hbar\omega}{c}\cos\theta$.

\textbf{3. Полный импульс} за время $\tau$ равен $pS\tau$:
\[pS\tau=\int_0^{\pi/2}2\frac{\hbar\omega}{c}\cos\theta\cdot nSc\tau\cos\theta\cdot\frac12\sin\theta\,\dd\theta
=nS\tau\hbar\omega\int_0^{\pi/2}\cos^2\theta\sin\theta\,\dd\theta=\frac{nS\tau\hbar\omega}{3}.\]
\answer{$p=\dfrac{n\hbar\omega}{3}=\dfrac u3$, где $u=n\hbar\omega$ — плотность энергии излучения.}
\note{для одноатомного газа аналогично $p=\frac23u$.}
\end{problem}

%-------------------------------------------------------------
\begin{problem}{8.12}{ex}{В сосуде при температуре $T$ находятся $N_1$ молекул массой $m_1$ и $N_2$ молекул массой $m_2$. Найти плотность
распределения $\tilde w(v)$ всех молекул по скоростям.}
\textbf{1.} Число молекул с модулем скорости в $(v,v+\dd v)$ складывается из молекул обоих сортов:
$\dd N=\dd N_1+\dd N_2=N_1w_1(v)\dd v+N_2w_2(v)\dd v$, где $w_i$ — распределение Максвелла с массой $m_i$.

\textbf{2.} С другой стороны, $\dd N=(N_1+N_2)\tilde w(v)\dd v$. Сравниваем.
\answer{$\tilde w(v)=\dfrac{N_1w_1(v)+N_2w_2(v)}{N_1+N_2}$, $w_i(v)=4\pi\Big(\dfrac{m_i}{2\pi kT}\Big)^{3/2}v^2e^{-m_iv^2/2kT}$.}
\note{если массы сильно различаются — у графика два максимума.}
\end{problem}

%-------------------------------------------------------------
\begin{problem}{8.19}{dz}{Найти скорость $\tilde v$ из условия, что среднее число молекул со скоростями меньше $\tilde v$ равно половине их общего числа.
Масса молекулы $m$, температура $T$.}
\textbf{Найти:} $\tilde v$ из условия $\int_0^{\tilde v}w(v)\dd v=\frac12$.

\textbf{1. Безразмерная переменная} $x=v/\vp$, $\vp=\sqrt{2kT/m}$: тогда $\frac{mv^2}{2kT}=x^2$ и
\[\int_0^{\tilde v}w(v)\dd v=\frac4{\sqrt\pi}\int_0^{x_0}x^2e^{-x^2}\dd x,\qquad x_0=\tilde v/\vp.\]

\textbf{2. По частям} (как в 8.4): $\int_0^{x_0}x^2e^{-x^2}\dd x=-\frac{x_0}2e^{-x_0^2}+\frac{\sqrt\pi}4\erf x_0$, поэтому условие:
\[\erf x_0-\frac2{\sqrt\pi}x_0e^{-x_0^2}=\frac12.\]

\textbf{3.} Уравнение трансцендентное, решаем численно (подбором по таблице $\erf$): $x_0\approx1{,}088$.
\answer{$\tilde v=x_0\sqrt{2kT/m}\approx1{,}088\sqrt{2kT/m}$.}
\note{$\vp<\tilde v<\vav<v_{\text{кв}}$: в единицах $\vp$ это $1<1{,}088<1{,}128<1{,}225$.}
\end{problem}

%-------------------------------------------------------------
\begin{problem}{8.20}{dz}{Вычислить средние значения проекции и модуля проекции скорости молекулы идеального газа на некоторое направление.
Масса молекулы $m$, температура $T$.}
\textbf{1. Среднее $v_x$.} $w(v_x)$ — чётная функция, а $v_x$ — нечётная, интеграл по симметричному промежутку равен нулю: $\avg{v_x}=0$.

\textbf{2. Среднее $|v_x|$.} Подынтегральная функция чётная — берём удвоенный интеграл по $v_x>0$:
\[\avg{|v_x|}=2\sqrt{\frac m{2\pi kT}}\int_0^\infty v_xe^{-mv_x^2/2kT}\dd v_x=2\sqrt{\frac m{2\pi kT}}\cdot\frac{kT}{m}=\sqrt{\frac{2kT}{\pi m}}\]
(использовали $\int_0^\infty xe^{-ax^2}\dd x=\frac1{2a}$, $a=\frac m{2kT}$).
\answer{$\avg{v_x}=0$,\; $\avg{|v_x|}=\sqrt{\dfrac{2kT}{\pi m}}=\dfrac{\vav}2$.}
\end{problem}

%-------------------------------------------------------------
\begin{problem}{8.21}{dz}{В сосуде при температуре $T$ находится $N$ молекул идеального газа. Найти число частиц, у которых одновременно проекция
скорости на ось $Oz$ лежит в $(v_z,v_z+\dd v_z)$, а модуль перпендикулярной составляющей скорости — в $(v_\perp,v_\perp+\dd v_\perp)$. Масса молекулы $m$.}
\textbf{1. Исходное распределение} по вектору скорости:
$\dd N=N\big(\frac m{2\pi kT}\big)^{3/2}e^{-m(v_x^2+v_y^2+v_z^2)/2kT}\dd v_x\dd v_y\dd v_z$.

\textbf{2. Цилиндрические координаты} в пространстве скоростей: $v_x=v_\perp\cos\varphi$, $v_y=v_\perp\sin\varphi$;
$v_x^2+v_y^2=v_\perp^2$, $\dd v_x\dd v_y=v_\perp\dd v_\perp\dd\varphi$.

\textbf{3. Направление $\vec v_\perp$ не важно} — интегрируем по $\varphi$ от 0 до $2\pi$, получаем множитель $2\pi$.
\answer{$\dd N=2\pi N\Big(\dfrac m{2\pi kT}\Big)^{3/2}v_\perp\exp\Big[-\dfrac{m(v_z^2+v_\perp^2)}{2kT}\Big]\dd v_z\,\dd v_\perp$.}
\end{problem}

%-------------------------------------------------------------
\begin{problem}{8.23}{dz}{Проекции скорости молекул идеального газа на оси статистически независимы и распределены по Гауссу. Найти совместную
$w(v_x,v)$ и условную $w(v_x\,|\,v)$ плотности вероятности. Масса молекулы $m$, температура $T$.}
\textbf{1. Условная плотность.} Если модуль $v$ задан, то направление скорости равновероятно по всем направлениям.
Для равномерного распределения на сфере $\cos\theta$ равномерно распределён на $[-1,1]$
(площадь пояса сферы $2\pi\sin\theta\,\dd\theta=2\pi\,|\dd\cos\theta|$). Так как $v_x=v\cos\theta$, то $v_x$ равномерно распределена на $[-v,v]$:
\[w(v_x\,|\,v)=\frac{1}{2v},\qquad|v_x|\le v.\]

\textbf{2. Совместная плотность} = плотность $v$ $\times$ условная: $w(v_x,v)=w(v)\cdot w(v_x\,|\,v)$, где $w(v)$ — распределение Максвелла:
\[w(v_x,v)=4\pi\Big(\frac m{2\pi kT}\Big)^{3/2}v^2e^{-mv^2/2kT}\cdot\frac1{2v}.\]
\answer{$w(v_x\,|\,v)=\dfrac1{2v}$,\; $w(v_x,v)=2\pi\Big(\dfrac m{2\pi kT}\Big)^{3/2}v\,e^{-mv^2/2kT}$\; (при $|v_x|\le v$, иначе 0).}
\end{problem}

\section*{§9. Распределение Больцмана}
\addcontentsline{toc}{section}{§9. Распределение Больцмана}

\begin{basics}
\sub{Распределение Больцмана}
\[n(\vec r)=n_0\,e^{-U(\vec r)/kT},\qquad
w(\vec r)=\frac{e^{-U(\vec r)/kT}}{\int_Ve^{-U/kT}\dd^3r},\qquad n(\vec r)=Nw(\vec r).\]
\begin{itemize}
\item $U(\vec r)$ — \emph{потенциальная энергия одной молекулы} во внешнем поле (Дж), а не потенциал поля.
Задана с точностью до константы (она сокращается). Если дана сила: $\vec F=-\nabla U$.
\item $n(\vec r)$ — концентрация (молекул на м$^3$); $n_0$ — концентрация там, где $U=0$.
\item $w(\vec r)$ — плотность вероятности найти \emph{данную} молекулу в точке $\vec r$: $w\,\dd^3r$ — вероятность
оказаться в объёме $\dd^3r$. Нормировка $\int_Vw\,\dd^3r=1$, т.е.\ $\int_Vn\,\dd^3r=N$ — отсюда находят $n_0$.
\item Применимость: идеальный газ (взаимодействие молекул мало), равновесие, $T$ одинакова везде.
\item Давление: $p(\vec r)=n(\vec r)kT$ — меняется так же, как $n$.
\end{itemize}

\sub{Не путать плотности}
\begin{tabular}{@{}lll@{}}
$w(\vec r)$ & на единицу объёма & $\propto e^{-U/kT}$\\
$w(z)$ & по высоте & $\propto S(z)e^{-U/kT}$ ($S$ — сечение)\\
$w(r)$ (шар) & по расстоянию от центра & $\propto4\pi r^2e^{-U/kT}$\\
$w(r)$ (цилиндр) & по расстоянию от оси & $\propto2\pi r\,e^{-U/kT}$
\end{tabular}
\\«Наиболее вероятное расстояние» — максимум $w(r)$ \emph{с множителем} $r^2$ или $r$.

\sub{Поле тяжести}
$U=mgz$ ($z$ — высота вверх):
\[n(z)=n(0)e^{-mgz/kT},\qquad p(z)=p(0)e^{-\mu gz/RT}\ \text{(барометрическая формула)}.\]
Полубесконечный цилиндр: $w(z)=\frac{mg}{kT}e^{-mgz/kT}$, $\avg z=\frac{kT}{mg}$, $\avg{z^2}=2\big(\frac{kT}{mg}\big)^2$, $\avg{mgz}=kT$.
\\Столб высоты $H$, $x=mgH/kT$: $\avg{mgz}=kT-\dfrac{mgH}{e^{x}-1}$.
\\Центр масс: $z_c=\dfrac{\sum m_i\avg{z_i}}{\sum m_i}$, для одного газа $z_c=\avg z$.

\sub{Вращающийся сосуд (центрифуга)}
В системе, вращающейся с $\omega$, на молекулу действует центробежная сила $m\omega^2r$, ей отвечает
$U=-\frac{m\omega^2r^2}{2}$ ($r$ — расстояние от оси):
$n(r)=n(0)\,e^{m\omega^2r^2/2kT}$ — газ прижимается к стенке, тяжёлые молекулы сильнее.

\sub{Несколько полей}
$U=U_1+U_2$ $\Rightarrow$ $e^{-U/kT}=e^{-U_1/kT}e^{-U_2/kT}$: координаты независимы, средние энергии складываются.
Полная энергия молекулы $\avg E=\frac32kT+\avg U$.
\end{basics}

%-------------------------------------------------------------
\begin{problem}{9.1}{exdz}{Найти среднее значение и относительную флуктуацию потенциальной энергии молекул идеального газа,
находящегося в прямом полубесконечном вертикальном цилиндре в однородном поле тяжести. Температура $T$, масса молекулы $m$.}
\textbf{Найти:} $\avg U$ и $\delta(U)=\sigma(U)/\avg U$, где $U=mgz$. Ось $z$ — от дна вверх, $z\in[0,\infty)$; $a=kT/mg$.

\textbf{1. Плотность вероятности по высоте.} По Больцману $w(z)=Ce^{-mgz/kT}=Ce^{-z/a}$ (сечение постоянное и сокращается).
Нормировка: $\int_0^\infty Ce^{-z/a}\dd z=Ca=1\Rightarrow w(z)=\frac1ae^{-z/a}$.

\textbf{2. Средние $z$ и $z^2$.} По формуле $\int_0^\infty z^ke^{-z/a}\dd z=k!\,a^{k+1}$:
\[\avg z=\frac1a\cdot a^2=a,\qquad \avg{z^2}=\frac1a\cdot2a^3=2a^2.\]

\textbf{3. Среднее и дисперсия $U=mgz$.} Множитель $mg$ выносится:
\[\avg U=mg\,\avg z=mga=kT,\qquad
\sigma^2(U)=\avg{U^2}-\avg U^2=(mg)^2\big(\avg{z^2}-\avg z^2\big)=(mg)^2a^2=(kT)^2.\]

\textbf{4. Относительная флуктуация:} $\delta(U)=\sigma(U)/\avg U=kT/kT=1$.
\answer{$\avg U=kT$, $\delta(U)=1$.}
\note{разброс энергии одной молекулы равен её среднему значению (100\%).}
\end{problem}

%-------------------------------------------------------------
\begin{problem}{9.2}{exdz}{Идеальный газ из $N$ молекул массой $m$ находится в вертикальном цилиндре высотой $H$ с площадью основания $\pi R^2$.
Температура $T$. Найти среднюю потенциальную энергию молекулы $\avg U$, высоту центра масс газа $z_c$ и во сколько раз
концентрация на высоте $z_c$ меньше, чем у дна: $\alpha=n(z_c)/n(0)$.}
Ось $z$ — от дна вверх. Обозначим $a=kT/mg$.

\textbf{1. Плотность вероятности по высоте.} По Больцману $w(z)=Ce^{-mgz/kT}=Ce^{-z/a}$ (сечение постоянное и сокращается).
Молекулы есть только при $0\le z\le H$, поэтому нормируем на этом отрезке:
\[\int_0^HCe^{-z/a}\dd z=Ca\big(1-e^{-H/a}\big)=1\ \Rightarrow\
w(z)=\frac{mg}{kT}\,\frac{e^{-mgz/kT}}{1-e^{-mgH/kT}}.\]

\textbf{2. Средняя высота.} $\avg z=\int_0^Hz\,w(z)\dd z$. Интеграл берём по частям:
\[\int_0^Hze^{-z/a}\dd z=\Big[-aze^{-z/a}\Big]_0^H+a\int_0^He^{-z/a}\dd z=-aHe^{-H/a}+a^2\big(1-e^{-H/a}\big).\]
Делим на нормировку $a(1-e^{-H/a})$:\quad $\avg z=a-\dfrac{H}{e^{H/a}-1}$.

\textbf{3. Ответы.} $\avg U=mg\avg z$; центр масс одинаковых молекул — их средняя высота, $z_c=\avg z$;
концентрация $n\propto e^{-mgz/kT}$, поэтому $\alpha=e^{-mgz_c/kT}$.
\answer{$\avg U=kT-\dfrac{mgH}{e^{mgH/kT}-1}$,\; $z_c=\dfrac{kT}{mg}-\dfrac{H}{e^{mgH/kT}-1}$,\; $\alpha=e^{-mgz_c/kT}$.}
\note{при $H\to\infty$ получается 9.1: $\avg U=kT$, $z_c=kT/mg$.}
\end{problem}

%-------------------------------------------------------------
\begin{problem}{9.3}{exdz}{Пылинки массой $m_0=10^{-21}$\,кг взвешены в воздухе при $T=300$\,К. Оценить толщину $\Delta z$ слоя воздуха у
поверхности Земли, в пределах которого концентрация пылинок отличается не более чем на 1\%.}
\textbf{Найти:} $\Delta z$, при котором $\Delta n/n=0{,}01$.

\textbf{1.} Пылинки ведут себя как молекулы газа в поле тяжести — распределение Больцмана:
$n(z)=n(0)\,e^{-m_0gz/kT}$.

\textbf{2.} Слой тонкий, поэтому изменение концентрации $\approx$ производная $\times$ толщина:
\[\Delta n\approx\Big|\frac{\dd n}{\dd z}\Big|\Delta z=\frac{m_0g}{kT}\,n\,\Delta z\quad\Rightarrow\quad
\frac{\Delta n}{n}=\frac{m_0g}{kT}\Delta z.\]

\textbf{3.} Выражаем $\Delta z$ и подставляем числа: $kT=1{,}38\cdot10^{-23}\cdot300=4{,}14\cdot10^{-21}$\,Дж,
$m_0g=9{,}8\cdot10^{-21}$\,Н:
\[\Delta z=\frac{kT}{m_0g}\cdot\frac{\Delta n}{n}=0{,}42\text{ м}\cdot0{,}01.\]
\answer{$\Delta z\approx4{,}2$ мм.}
\end{problem}

%-------------------------------------------------------------
\begin{problem}{9.4}{exdz}{В центрифуге радиуса $R$ и высоты $H$, вращающейся с угловой скоростью $\omega$, разделяется смесь двух газов
с массами молекул $m_1$ и $m_2$. Найти коэффициент разделения $q=\dfrac{n_1(R)/n_1(0)}{n_2(R)/n_2(0)}$, где $n_{1,2}(r)$ — концентрации
газов на расстоянии $r$ от оси. Силу тяжести не учитывать, температура $T$.}
\textbf{1. Потенциальная энергия.} Во вращающейся системе на молекулу действует центробежная сила $F=m_i\omega^2r$ (от оси).
Ей соответствует $U_i=-\int F\,\dd r=-\dfrac{m_i\omega^2r^2}{2}$.

\textbf{2. Больцман для каждого газа отдельно} (газы идеальные, не влияют друг на друга):
\[n_i(r)=n_i(0)\,e^{-U_i/kT}=n_i(0)\,e^{m_i\omega^2r^2/2kT}\ \Rightarrow\ \frac{n_i(R)}{n_i(0)}=e^{m_i\omega^2R^2/2kT}.\]

\textbf{3. Делим} выражение для газа 1 на выражение для газа 2.
\answer{$q=\exp\Big[\dfrac{(m_1-m_2)\omega^2R^2}{2kT}\Big]$.}
\note{UF$_6$ с $^{235}$U и $^{238}$U, $R=1$\,м, $T=300$\,К, $\omega=1300$\,с$^{-1}$: $|m_1-m_2|=3\cdot1{,}66\cdot10^{-27}$\,кг, показатель $\approx1{,}0$;
в книге $q\approx2{,}7$ — это если тяжёлый изотоп поставить в числитель.}
\end{problem}

%-------------------------------------------------------------
\begin{problem}{9.5}{ex}{Смесь двух идеальных газов: $N_1$ частиц массой $m_1$ и $N_2$ частиц массой $m_2$ — находится в прямом
цилиндрическом сосуде (высоту считаем очень большой). Найти положение центра масс смеси, если температура $T$, $\vec g$ направлено вдоль оси.}
\textbf{1. Центр масс} — по определению среднее высот, взвешенное по массам:
\[z_c=\frac{N_1m_1\avg{z_1}+N_2m_2\avg{z_2}}{N_1m_1+N_2m_2}.\]

\textbf{2. Средняя высота молекулы каждого газа} — из 9.1 (газы независимы): $\avg{z_i}=\dfrac{kT}{m_ig}$, поэтому $m_i\avg{z_i}=\dfrac{kT}{g}$ для обоих.
\answer{$z_c=\dfrac{(N_1+N_2)\,kT}{g\,(N_1m_1+N_2m_2)}$.}
\note{при конечной высоте $H$ брать $\avg{z_i}$ из 9.2.}
\end{problem}

%-------------------------------------------------------------
\begin{problem}{9.8}{dz}{Идеальный газ из $N$ частиц находится в сосуде в виде перевёрнутого прямого конуса (вершина внизу, на поверхности Земли),
ось параллельна $\vec g$, угол при вершине $2\beta$, высота бесконечно большая. Найти концентрацию $n_0$ молекул у вершины. Температура $T$, масса молекулы $m$.}
\textbf{1. Больцман:} $n(z)=n_0e^{-mgz/kT}$, $z$ — высота над вершиной.

\textbf{2. Объём тонкого слоя.} На высоте $z$ сечение — круг радиуса $z\tan\beta$, площадь $S(z)=\pi z^2\tan^2\beta$, объём слоя $S(z)\dd z$.

\textbf{3. Условие нормировки:} всех молекул $N$:
\[N=\int_0^\infty n(z)S(z)\dd z=n_0\pi\tan^2\beta\int_0^\infty z^2e^{-mgz/kT}\dd z=n_0\pi\tan^2\beta\cdot2\Big(\frac{kT}{mg}\Big)^3\]
(использовали $\int_0^\infty z^2e^{-z/a}\dd z=2a^3$, $a=kT/mg$). Выражаем $n_0$.
\answer{$n_0=\dfrac{N(mg)^3}{2\pi\tan^2\beta\,(kT)^3}$.}
\note{так же 9.7 (цилиндр с площадью дна $S$): $N=n_0S\frac{kT}{mg}\Rightarrow n_0=\frac{Nmg}{SkT}$.}
\end{problem}

%-------------------------------------------------------------
\begin{problem}{9.12}{dz}{Отрицательные ионы массой $m$ и зарядом $q$ находятся при температуре $T$ между пластинами плоского конденсатора.
Поле $\vec E$ направлено туда же, куда $\vec g$ (вниз). Во сколько раз концентрация ионов у верхней пластины отличается от
концентрации у нижней, если расстояние между пластинами $d$?}
\textbf{Найти:} $n_{\text{верх}}/n_{\text{низ}}$. Ось $z$ — вверх, $z=0$ у нижней пластины.

\textbf{1. Силы.} Тяжесть $mg$ — вниз. Ион отрицательный, поэтому электрическая сила $|q|E$ направлена \emph{против} $\vec E$, т.е.\ вверх.

\textbf{2. Потенциальная энергия} (сила вниз $\Rightarrow U$ растёт с высотой):
$U(z)=mgz-|q|Ez=(mg-|q|E)z$.

\textbf{3. Больцман:} $\dfrac{n(d)}{n(0)}=e^{-[U(d)-U(0)]/kT}=e^{-(mg-|q|E)d/kT}$.
\answer{$\dfrac{n_{\text{верх}}}{n_{\text{низ}}}=\exp\Big[\dfrac{d\,(|q|E-mg)}{kT}\Big]$.}
\end{problem}

%-------------------------------------------------------------
\begin{problem}{9.13}{dz}{Идеальный газ находится в поле центральной силы с потенциальной энергией $U(r)=\alpha r^2$ ($\alpha>0$, $r$ — расстояние до центра).
Температура $T$, число частиц $N$. Найти $n(r)$ и наиболее вероятное расстояние $r_0$ молекул от центра. Найти плотность вероятности $w(U)$
и наиболее вероятное значение потенциальной энергии $\tilde U$.}
\textbf{1. Концентрация.} По Больцману $n(r)=n_0e^{-\alpha r^2/kT}$. Константу $n_0$ — из нормировки по всему пространству (слой $4\pi r^2\dd r$):
\[N=n_0\int_0^\infty4\pi r^2e^{-\alpha r^2/kT}\dd r=n_0\cdot4\pi\cdot\frac{\sqrt\pi}{4}\Big(\frac{kT}{\alpha}\Big)^{3/2}=n_0\Big(\frac{\pi kT}{\alpha}\Big)^{3/2}\]
(гауссов интеграл $I_2=\frac{\sqrt\pi}{4a^{3/2}}$ с $a=\alpha/kT$).

\textbf{2. Наиболее вероятное расстояние.} Вероятность оказаться в слое $(r,r+\dd r)$ пропорциональна \emph{объёму слоя}:
$w(r)\propto r^2e^{-\alpha r^2/kT}$. Максимум: $\big(2r-\frac{2\alpha r^3}{kT}\big)e^{\dots}=0\Rightarrow r_0=\sqrt{kT/\alpha}$.

\textbf{3. Распределение по $U$.} Замена $U=\alpha r^2$: $r=\sqrt{U/\alpha}$, $\dd r=\frac{\dd U}{2\sqrt{\alpha U}}$:
\[w(U)\propto r^2e^{-U/kT}\frac{\dd r}{\dd U}\propto\frac U\alpha\cdot\frac{1}{\sqrt{\alpha U}}e^{-U/kT}\propto\sqrt U\,e^{-U/kT}.\]
Это та же форма, что $w(E)$ в 8.2, поэтому нормировочная константа та же. Максимум: $\frac{\dd}{\dd U}\big(\sqrt Ue^{-U/kT}\big)=0\Rightarrow\tilde U=\frac{kT}2$.
\answer{$n(r)=N\Big(\dfrac{\alpha}{\pi kT}\Big)^{3/2}e^{-\alpha r^2/kT}$,\; $r_0=\sqrt{kT/\alpha}$,\;
$w(U)=\dfrac{2}{\sqrt\pi}(kT)^{-3/2}\sqrt U\,e^{-U/kT}$,\; $\tilde U=\dfrac{kT}2$.}
\note{$\tilde U\ne\alpha r_0^2=kT$: наиболее вероятные значения не пересчитываются через замену.}
\end{problem}

%-------------------------------------------------------------
\begin{problem}{9.14}{dz}{Газ находится в полубесконечном прямом цилиндре в однородном поле тяжести. Найти долю молекул, потенциальная
энергия которых больше их средней кинетической энергии.}
\textbf{Найти:} долю молекул с $mgz>\avg{E_{\text{кин}}}=\frac32kT$, т.е.\ выше высоты $z_0=\frac{3kT}{2mg}$.

\textbf{1.} Плотность по высоте (9.1): $w(z)=\frac{mg}{kT}e^{-mgz/kT}$.

\textbf{2.} Доля выше $z_0$:
$\displaystyle P=\int_{z_0}^\infty w(z)\dd z=e^{-mgz_0/kT}=e^{-3/2}$.
\answer{$e^{-1{,}5}\approx0{,}22$.}
\end{problem}

%-------------------------------------------------------------
\begin{problem}{9.15}{dz}{В центрифуге радиуса $R$ и высоты $H$, вращающейся с угловой скоростью $\omega$, при температуре $T$ находится $N$ молекул
идеального газа массой $m$. Найти среднее значение полной энергии молекулы. Ускорение $\vec g$ направлено вдоль оси.}
\textbf{1. Из чего состоит энергия.} $E=E_{\text{кин}}+U_{\text{ц}}+U_g$, где $U_{\text{ц}}=-\frac{m\omega^2r^2}2$ (центробежная), $U_g=mgz$.
Среднее суммы = сумма средних. Кроме того, $e^{-(U_{\text{ц}}+U_g)/kT}=e^{-U_{\text{ц}}/kT}e^{-U_g/kT}$ — по $r$ и по $z$ распределения независимы,
каждое слагаемое считаем отдельно. $\avg{E_{\text{кин}}}=\frac32kT$.

\textbf{2. Тяжесть} — ответ 9.2: $\avg{U_g}=kT-\dfrac{mgH}{e^{mgH/kT}-1}$.

\textbf{3. Центробежная часть.} По радиусу: $w(r)\propto2\pi r\,e^{m\omega^2r^2/2kT}$ (площадь кольца $2\pi r\,\dd r$).
Удобная переменная $s=\frac{m\omega^2r^2}{2kT}\in[0,b]$, $b=\frac{m\omega^2R^2}{2kT}$; так как $\dd s\propto r\,\dd r$, то $w(s)=\frac{e^s}{e^b-1}$. Тогда
\[\avg s=\frac{\int_0^bse^s\dd s}{e^b-1}=\frac{(b-1)e^b+1}{e^b-1}=\frac{b\,e^b}{e^b-1}-1,\qquad
\avg{U_{\text{ц}}}=-kT\avg s=kT-\frac{m\omega^2R^2}{2(1-e^{-b})}\]
(это ответ 9.9).

\textbf{4. Складываем:} $\frac32kT+kT+kT=\frac72kT$ плюс две поправки.
\answer{$\avg E=\dfrac{7kT}2-\dfrac{m\omega^2R^2}{2\big[1-e^{-m\omega^2R^2/2kT}\big]}-\dfrac{mgH}{e^{mgH/kT}-1}$.}
\end{problem}

%-------------------------------------------------------------
\begin{problem}{9.16}{dz}{Два прямых цилиндрических сосуда высотой $H$ каждый стоят друг на друге и содержат по одному молю идеального газа при
температуре $T$. На какой высоте от поверхности Земли находится центр масс системы? Масса молекулы $m$.}
\textbf{1. Нижний сосуд} ($0\le z\le H$) — ответ 9.2: $z_1=\dfrac{kT}{mg}-\dfrac{H}{e^{x}-1}$, $x=\dfrac{mgH}{kT}$.

\textbf{2. Верхний сосуд} — то же распределение, но отсчитанное от его дна на высоте $H$: $z_2=z_1+H$.

\textbf{3. Массы газов одинаковы} (по молю), поэтому центр масс — посередине: $z_c=\frac{z_1+z_2}{2}=z_1+\frac H2$.
\answer{$z_c=\dfrac{kT}{mg}+\dfrac H2-\dfrac{H}{e^{mgH/kT}-1}$.}
\end{problem}

\end{document}
